\documentclass[11pt,a4paper]{article}
\usepackage[T1]{fontenc}
\usepackage[english,french]{babel}
\RequirePackage{euler,eulervm}
\usepackage{libertine}
\RequirePackage[italic]{mathastext}
\usepackage{graphicx}
\usepackage{amssymb,amsthm,amsmath}
\usepackage{multicol}
\usepackage{booktabs}
\usepackage{enumitem}
\usepackage{incgraph,tikz}
\usepackage{lscape}
\usepackage{caption}
\usepackage{calc}
\usepackage{colortbl}
\usepackage{natbib}
\usepackage{hyperref}
\usepackage{forest}
\usepackage{adjustbox}
\usepackage{setspace}
\usepackage[dvipsnames]{xcolor}
\usepackage{ragged2e}
\usepackage[skip=\bigskipamount,indent]{parskip}

\usetikzlibrary{arrows.meta}
\hypersetup{
	colorlinks = true,
	linkcolor= blue,
	urlcolor= blue,
	citecolor = blue}

\usepackage{multirow}
\usepackage{makecell}
\usepackage{dcolumn}
\usepackage[table]{xcolor}% http://ctan.org/pkg/xcolor
\usepackage{overpic}
\usepackage{pdflscape} % Add to your preamble
\usepackage{lscape}
\usepackage{afterpage}
\usepackage[usegeometry,DIV=15]{typearea}
\usepackage[scale=.85]{geometry}
\newcommand{\trj}{\mathit{trj}}
\newcommand{\agngvt}{\mathit{agngvt}}
\newcommand{\ag}{\mathit{ag}}
\newcommand{\agj}{\mathit{agj}}
\newcommand{\agd}{\mathit{agd}}





\newcommand{\agpub}{\mathit{agpub}}
\newcommand{\vel}{\mathit{vel}}



\newcommand{\secpub}{\mathit{secpub}}
\newcommand{\ntyh}{\mathit{ntyh}}
\newcommand{\ntyf}{\mathit{ntyf}}
\newcommand{\gvt}{\mathsf{gvt}}
\newcommand{\hur}{\mathsf{hur}}
\newcommand{\hru}{\mathsf{hru}}
\newcommand{\firm}{\mathsf{firm}}
\newcommand{\ngo}{\mathsf{ngo}}
\newcommand{\row}{\mathsf{row}}

\newcommand{\NTX}{\mathit{NTX}}

\newcommand{\YDF}{\mathit{YDF}}
\newcommand{\YFTR}{\mathit{YFTRA}}
\newcommand{\npbns}{\mathit{npbns}}
\newcommand{\aij}{\mathit{aij}}
\newcommand{\mij}{\mathit{mij}}
\newcommand{\bns}{\mathit{bns}}
\newcommand{\NDDG}{\mathit{NDDG}}
\newcommand{\CAB}{\mathit{CAB}}
\newcommand{\CI}{\mathit{CI}}
\newcommand{\CG}{\mathit{CG}}
\newcommand{\OTHNTR}{\mathit{OTHNTR}}
\newcommand{\CH}{\mathit{CH}}
\newcommand{\CNG}{\mathit{CNG}}
\newcommand{\CTH}{\mathit{CTH}}

\newcommand{\DD}{\mathit{DD}}

\newcommand{\DX}{\mathit{DX}}
\newcommand{\DI}{\mathit{DI}}

\newcommand{\ontr}{\mathsf{ontr}}
\newcommand{\sntx}{\mathit{sntx}}
\newcommand{\DIT}{\mathit{DIT}}
\newcommand{\DS}{\mathit{DS}}

\newcommand{\DEFG}{\mathit{DEFG}}
\newcommand{\EX}{\mathit{EX}}
\newcommand{\FBCF}{\mathit{FBCF}}
\newcommand{\BRG}{\mathit{BRG}}
\newcommand{\IM}{\mathit{IM}}
\newcommand{\TOTG}{\mathit{TOTG}}
\newcommand{\NTXHH}{\mathit{NTXHH}}
\newcommand{\INV}{\mathit{INV}}
\newcommand{\IT}{\mathit{IT}}
\newcommand{\io}{\mathit{io}}
\newcommand{\MGTRP}{\mathit{MGTRP}}
\newcommand{\MGCOM}{\mathit{MGCOM}}
\newcommand{\MRGTRP}{\mathit{MRGTRP}}
\newcommand{\MRGCOM}{\mathit{MRGCOM}}
\newcommand{\DTH}{\mathit{DTH}}
\newcommand{\DTF}{\mathit{DTF}}
\newcommand{\KD}{\mathit{KD}}
\newcommand{\KS}{\mathit{KS}}
\newcommand{\LD}{\mathit{LD}}
\newcommand{\LEON}{\mathit{LEON}}
\newcommand{\LEONF}{\mathit{LEONF}}
\newcommand{\LS}{\mathit{LS}}



\newcommand{\NTR}{\mathit{NTR}}
\newcommand{\ntr}{\mathit{ntr}}

% Epargnes
\newcommand{\othntr}{\mathsf{othntr}}
\newcommand{\agr}{\mathsf{agr}}
\newcommand{\food}{\mathsf{food}}
\newcommand{\othind}{\mathsf{othind}}
\newcommand{\ser}{\mathsf{ser}}

\newcommand{\DINVG}{\mathit{DINVG}}
\newcommand{\FINVG}{\mathit{FINVG}}
\newcommand{\SNG}{\mathit{SNG}}
\newcommand{\SNGO}{\mathit{SNGO}}
\newcommand{\snb}{\mathit{snb}}
\newcommand{\SF}{\mathit{SF}}
\newcommand{\SG}{\mathit{SG}}
\newcommand{\SH}{\mathit{SH}}
\newcommand{\TD}{\mathit{TD}}
\newcommand{\TI}{\mathit{TI}}
\newcommand{\TIF}{\mathit{TIF}}
\newcommand{\TIH}{\mathit{TIH}}
\newcommand{\TIM}{\mathit{TIM}}
\newcommand{\TIMF}{\mathit{TIMF}}
\newcommand{\PE}{\mathit{PE}}
\newcommand{\ONTR}{\mathit{ONTR}}
\newcommand{\PC}{\mathit{PC}}
\newcommand{\PCI}{\mathit{PCI}}
\newcommand{\PLD}{\mathit{PLD}}
\newcommand{\PD}{\mathit{PD}}
\newcommand{\PPD}{\mathit{PPD}}
\newcommand{\PDX}{\mathit{PDX}}
\newcommand{\PINDEX}{\mathit{PINDEX}}
\newcommand{\PIXPIB}{\mathit{PIXPIB}}
\newcommand{\PIXCON}{\mathit{PIXCON}}
\newcommand{\PIXINV}{\mathit{PIXINV}}
\newcommand{\PL}{\mathit{PL}}
\newcommand{\PM}{\mathit{PM}}
\newcommand{\TRA}{\mathit{TRA}}
\newcommand{\PPM}{\mathit{PPM}}
\newcommand{\PWE}{\mathit{PWE}}
\newcommand{\PWM}{\mathit{PWM}}
\newcommand{\pub}{\mathsf{pub}}
\newcommand{\PVA}{\mathit{PVA}}
%\newcommand{\ser}{\mathit{ser}}
\newcommand{\TCM}{\mathit{TCM}}
\newcommand{\te}{\mathit{te}}
\newcommand{\TG}{\mathit{TG}}
\newcommand{\TR}{\mathit{TRA}}
\newcommand{\TRG}{\mathit{TRG}}
\newcommand{\TRH}{\mathit{TRH}}
\newcommand{\TRNG}{\mathit{TRNG}}
\newcommand{\TRF}{\mathit{TRF}}
\newcommand{\syng}{\mathit{syng}}
\newcommand{\syf}{\mathit{syf}}
\newcommand{\petr}{\mathsf{petr}}
\newcommand{\syfi}{\mathit{syfi}}
\newcommand{\syfm}{\mathit{syfm}}
\newcommand{\syh}{\mathit{syh}}
\newcommand{\tr}{\mathit{tr}}
\newcommand{\ytrf}{\mathit{ytrf}}
\newcommand{\tx}{\mathit{tx}}
\newcommand{\tm}{\mathit{tm}}
\newcommand{\tmgc}{\mathit{tmgc}}
\newcommand{\tmgt}{\mathit{tmgt}}
\newcommand{\tyh}{\mathit{tyh}}
\newcommand{\tyf}{\mathit{tyf}}
\newcommand{\VA}{\mathit{VA}}
\newcommand{\VSTK}{\mathit{VSTK}}
\newcommand{\WLG}{\mathit{WLG}}
\newcommand{\ROW}{\mathit{ROW}}
\newcommand{\XS}{\mathit{XS}}
\newcommand{\NTXF}{\mathit{NTXF}}
\newcommand{\NTXH}{\mathit{NTXH}}
\newcommand{\XD}{\mathit{XD}}
\newcommand{\XSP}{\mathit{XSP}}
\newcommand{\XSPT}{\mathit{XSPT}}
\newcommand{\TDH}{\mathit{TDH}}
\newcommand{\TDF}{\mathit{TDF}}
\newcommand{\YDH}{\mathit{YDH}}
\newcommand{\YH}{\mathit{YH}}
\newcommand{\YHK}{\mathit{YHK}}
\newcommand{\YHL}{\mathit{YHL}}
\newcommand{\YHTR}{\mathit{YHTRA}}
\newcommand{\YF}{\mathit{YF}}
\newcommand{\YFK}{\mathit{YFK}}
\newcommand{\YG}{\mathit{YG}}
\newcommand{\YGT}{\mathit{YGT}}
\newcommand{\YGTX}{\mathit{YGTX}}
\newcommand{\YGNTX}{\mathit{YGNTX}}
\newcommand{\YGTR}{\mathit{YGTRA}}
\newcommand{\YNG}{\mathit{YNG}}
\newcommand{\YNGO}{\mathit{YNGO}}
\newcommand{\YROW}{\mathit{YROW}}
\newcommand{\YROWK}{\mathit{YROWK}}
\newcommand{\YROWL}{\mathit{YROWL}}
\newcommand{\SROW}{\mathit{SROW}}
\newcommand{\redparagraph}[1]{\paragraph*{\large\textcolor{Fuchsia}{#1}}}
\newcommand{\PIB}{\mathit{PIB}}


\newcommand{\INVG}{\mathit{INVG}}
\newcommand{\KTRG}{\mathit{KTRG}}
\newcommand{\NFDG}{\mathit{NFDG}}
\newcommand{\DCG}{\mathit{DCG}}
\newcommand{\BG}{\mathit{BG}}
\newcommand{\AMORTG}{\mathit{AMORTG}}
\newcommand{\SP}{\mathit{SP}}


\newcommand{\RES}{\mathit{RES}}
\newcommand{\OTHKTR}{\mathit{OTHKTR}}
\newcommand{\FDI}{\mathit{FDI}}
\newcommand{\NFDP}{\mathit{NFDP}}

\newcommand{\MS}{\mathit{MS}}
\newcommand{\MD}{\mathit{MD}}
\newcommand{\DCP}{\mathit{DCP}}
\newcommand{\INVP}{\mathit{INVP}}

\newcommand{\ER}{\mathit{ER}}
\newcommand{\DEBTG}{\mathit{DEBTG}}

\everymath\expandafter{\the\everymath\displaystyle}
\usepackage{tikz}
\usetikzlibrary{shapes.geometric, arrows.meta, positioning, matrix, calc,decorations.pathreplacing, fit}


\usepackage{orcidlink}

\makeatletter
\newcolumntype{B}[3]{>{\boldmath\DC@{#1}{#2}{#3}}c<{\DC@end}}
\makeatother
\newcommand{\gras}[1]{\multicolumn{1}{B{.}{.}{5.1}|}{#1}}
\newcommand{\grasl}[1]{\multicolumn{1}{@{\vrule width 1.5pt}B{.}{.}{5.1}|}{#1}}
\newcommand{\grascol}[1]{\multicolumn{1}{B{.}{.}{5.2}@{\vrule width 1.5pt}}{#1}}



\title{EXTER-MADA : Un modèle d'équilibre général calculable pour l'analyse des politiques économiques à Madagascar}
\author{Franck M.~Ramaharo\orcidlink{0000-0003-0340-170X}\\
	\small Service de la Modélisation Économique\\
	\small Ministère de l'Économie et des Finances\\ 
	\small 101 Antananarivo, Madagascar\\
	\small\href{mailto:franck.ramaharo@gmail.com}{\sffamily franck.ramaharo@gmail.com}\\
}

\date{\today}
\begin{document}
	\maketitle
	
	
	\begin{abstract}
		Ce document présente EXTER-MADA, un modèle d’équilibre général calculable (MEGC) statique adapté à l’économie malgache et calibré sur une Matrice de Comptabilité Sociale (MCS) de 2019. Le modèle est dérivé des modèles EXTER de \cite{Decaluweetal2001} et PEP-1-1 de \cite{Decaluweetal2013}, tout en retenant les spécificités d'anciens modèles macroéconomiques nationaux.  
		
		\bigskip
		
		\textbf{Mots-clés :} Madagascar, modèle d’équilibre général calculable, MEGC, matrice de comptabilité sociale, MCS, modélisation macroéconomique
	\end{abstract}	
	
	
	\selectlanguage{english} 
	\begin{abstract}
		
		This document presents EXTER-MADA, a static computable general equilibrium (CGE) model tailored to the Malagasy economy and calibrated on a 2019 Social Accounting Matrix (SAM). The model is derived from the EXTER model by \cite{Decaluweetal2001} and the PEP-1-1 model by \cite{Decaluweetal2013}, while incorporating specific features from earlier national macroeconomic models.
		
		\bigskip
		
		\textbf{Keywords:} Madagascar,  computable general equilibrium model, CGE, social accounting matrix, SAM, macroeconomic modeling
	\end{abstract}
	\selectlanguage{french} 
	
	\vfill
	\begin{center}
		\fbox{\noindent\begin{minipage}[c]{\textwidth}
				\textbf{Note de mise en garde} : Ce document est un travail en cours (\textit{work in progress}). Les développements théoriques, la modélisation et les résultats empiriques sont présentés à titre provisoire et sont appelés à évoluer. Afin d'en assurer la transparence et la réplicabilité, et suivant le principe de la science ouverte, l'ensemble du projet est hébergé en accès ouvert sur le dépôt OSF (Open Science Framework) suivant : \href{https://doi.org/10.17605/OSF.IO/AS3BH}{\sffamily https://doi.org/10.17605/OSF.IO/AS3BH}.\\
				La Matrice de Comptabilité Sociale (MCS) sous-jacente au modèle MEGC présenté dans ce document repose sur une estimation propre réalisée par l'auteur ; les notes méthodologiques d'accompagnement ainsi que les révisions de la base de données seront versées de manière continue dans le dépôt. Cette démarche vise également à encourager les observations, réflexions et suggestions des pairs destinées à amender les versions futures.
				\bigskip
				
				\textbf{Clause de non-responsabilité (\textit{Disclaimer})} : Les analyses, résultats, interprétations et conclusions exprimés dans ce document relèvent de la seule responsabilité de l'auteur. Ils ne reflètent pas nécessairement la position officielle ni les orientations stratégiques des institutions ou ministères auxquels l'auteur est affilié.
		\end{minipage}}	
	\end{center}
	\newpage
	\onehalfspacing
	\section{Introduction}	
	Le modèle EXTER-MADA présenté ici est un modèle d’équilibre général calculable (MEGC) statique appliqué à l'économie malgache. Sa structure théorique constitue une version intermédiaire entre le modèle de référence EXTER \citep{Decaluweetal2001} et le modèle PEP-1-1 \citep{Decaluweetal2013}. Sa conception intègre les enseignements tirés des modélisations antérieures réalisées pour Madagascar, notamment le modèle MADCGE initialement développé au sein de l'Institut National de la Statistique \citep{Dissou1998, Dissou2001,Dissouetal1998,Doroshetal2003}. Il retient également les structures et les règles quasi comptables de modèles tels que MAROA, MA2P et VENOM, lesquels ont été développés au sein du ministère en charge de l'Économie.
	
	Les hypothèses de la modélisation:
	\begin{itemize}
		\item Tous les agents évoluent dans un environnement de concurrence pure et parfaite. Les producteurs maximisent leur profit et les ménages leur utilité, sous l'hypothèse qu'aucun d'entre eux ne peut influencer les prix, qu'ils considèrent comme des données exogènes.
		
		
		\item À court-moyen terme, le capital est supposé spécifique à chaque secteur d'activité (\ref{it:KS}), limitant la réallocation instantanée des capacités de production. À l'inverse, le travail est considéré comme parfaitement mobile entre les secteurs (\ref{it:LS}).
		
		
		\item L'État est modélisé à travers le compte de l'Administration Centrale, qui collecte les recettes fiscales et non fiscales, et assure la fourniture des services publics.
		
		
		\item  Les secteurs produisant des biens échangeables arbitrent entre le marché domestique et le marché d'exportation. En revanche, les secteurs de biens non échangeables (administration publics, éducation publique, santé publique) produisent exclusivement pour le marché intérieur.
	\end{itemize} 
	%
	
	%
	\section{La Matrice de Comptabilité Sociale}
	La Matrice de Comptabilité Sociale (MCS) élaborée dans le cadre de ce modèle comprend les comptes suivants :
	\begin{itemize}
		\item[$\bullet$] \textbf{2 comptes de facteurs de production} : Travail et Capital ;
		\item[$\bullet$] \textbf{7 comptes d’agents économiques} : Ménages urbains, Ménages ruraux, Administration centrale, Autres Administrations Publiques (APU) et Institutions Sans But Lucratif au Service des Ménages (ISBLSM), Sociétés et Reste du monde ;
		\item[$\bullet$] \textbf{4 comptes d'impôts, taxes et recettes publiques} : Impôts directs, Taxes sur les biens et services, Taxes sur le commerce international et Ressources non fiscales ;
		\item[$\bullet$] \textbf{6 comptes de branches d’activité} : Agriculture, ``Alimentaire, Boisson, Tabac'', Autres Industries, Services marchands, Administration centrale et Autres services publics et services non marchands ;
		\item[$\bullet$] \textbf{5 comptes de produits domestiques} (vendus sur le marché local) : Agriculture, ``Alimentaire, Boisson, Tabac'', Autres Industries, Produits pétroliers, Service marchand, Administration publique et Autres services publics et services non marchands ;
		\item[$\bullet$] \textbf{5 comptes de produits d’exportation} (vendus sur le marché extérieur) : Agriculture, ``Alimentaire, Boisson, Tabac'', Autres Industries, Produits pétroliers et Services marchands;
		\item[$\bullet$] \textbf{2 comptes des marges} : Marge de commercialisation et Marge de transport ;
		\item[$\bullet$] \textbf{2 comptes d’accumulation} : Formation Brute de Capital Fixe (FBCF) et Variation de stock.
	\end{itemize}
	
	Les versions agrégée (MACROSAM) et désagrégée (MIDISAM) de la matrice de comptabilité sociale sont données dans les \textsc{Tableaux} \ref{tab:MACROSAM} et \ref{tab:MIDISAM}. Notons que la MACROSAM présentée dans le \textsc{Tableau}~\ref{tab:MACROSAM} est construite de manière similaire à celle de \cite{Ramaharo2022}. À partir d'un certain niveau d'agrégation, les résultats sont par construction identiques. Les  divergences statistiques peuvent être expliquées par la mise à jour des données de base ainsi que par quelques ajustements dans la définition des agrégats.
	
	
	\begin{table}[!htbp]
		\renewcommand{\arraystretch}{1.2} % Scales row height (1.0 is default)
		\centering
		\caption{Version aggrégée  de la Matrice de Comptabilité Sociale de Madagascar (MACROSAM). Année 2019. En Milliards d'Ariary.}
		\makebox[0pt]{\resizebox{1.1\columnwidth}{!}{%
				\begin{tabular}{|r|D{.}{.}{5.1}|D{.}{.}{5.1}|D{.}{.}{5.1}|D{.}{.}{5.1}|D{.}{.}{5.1}|D{.}{.}{5.1}|D{.}{.}{5.1}|D{.}{.}{5.1}|D{.}{.}{5.1}|D{.}{.}{5.1}|D{.}{.}{5.1}|D{.}{.}{5.1}|D{.}{.}{5.1}|}
					\cline{2-13}    \multicolumn{1}{r|}{} & \textbf{\sffamily1} & \textbf{\sffamily2} & \textbf{\sffamily3} & \textbf{\sffamily4} & \textbf{\sffamily5} & \textbf{\sffamily6} & \textbf{\sffamily7} & \textbf{\sffamily8} & \textbf{\sffamily9} & \textbf{\sffamily10} & \textbf{\sffamily11} & \textbf{\sffamily12} \\
					\hline
					\sffamily\textbf{1} &    &    &    &    &    & 221.6 & 47638.3 &    &    &    &    & \gras{47859.9} \\
					\hline
					\textbf{2} & 27522.1 &    & 677.5 & 13786.6 &    & 1266.8 &    &    &    &    &    & \gras{43253.0} \\
					\hline
					\sffamily\textbf{3} &    &  & 754.8 & 4331.1 & 5959.4 & 432.0 &    &    &    &    &    & \gras{11477.3} \\
					\hline
					\sffamily\textbf{4} & 18617.2 &    & 1292.5 &    &    & 1189.7 &    &    &    &    &    & \gras{21099.4} \\
					\hline
					\sffamily\textbf{5} &    & 1004.5 &    & 1557.9 &    &    &    & 3396.9 &    &    &    & \gras{5959.4} \\
					\hline
					\sffamily\textbf{6} & 1720.6 &    & 106.8 &    &    &    &    & 17442.3 &    &    &    & \gras{19269.7} \\
					\hline
					\sffamily\textbf{7} &    &    &    &    &    &    &    & 65405.2 & 14228.0 &    &    & \gras{79633.1} \\
					\hline
					\sffamily\textbf{8} &    & 34684.7 & 7709.4 &    &    &    & 31994.8 &    & 278.5 & 8531.2 & 11576.9 & \gras{94775.6} \\
					\hline
					\sffamily\textbf{9} &    &    &    &    &    & 14506.5 &    &    &    &    &    & \gras{14506.5} \\
					\hline
					\textbf{10} &    &    &    &    &    &    &    & 8531.2 &    &    &    & \gras{8531.2} \\
					\hline
					\sffamily\textbf{11} &    & 7563.8 & 936.1 & 1423.8 &    & 1653.2 &    &    &    &    & 733.0 & \gras{12309.8} \\
					\hline
					\sffamily\textbf{12} & \gras{47859.9} & \gras{43253.0} & \gras{11477.3} & \gras{21099.4} & \gras{5959.4} & \gras{19269.7} & \gras{79633.1} & \gras{94775.6} & \gras{14506.5} & \gras{8531.2} & \gras{12309.8} &  \\
					\hline
				\end{tabular}%
		}}
		\begin{multicols}{3}
			\begin{enumerate}[font=\bfseries\sffamily]\sffamily\small
				\item \textbf{Facteurs}
				\item \textbf{Ménages}
				\item \textbf{Administration\\ publique et ISBLSM}
				\item \textbf{Sociétés}
				\item \textbf{Prélèvements fiscaux et non fiscaux}
				\item \textbf{Reste du monde}
				\item \textbf{Activités}
				\item \textbf{Produits domestiques}
				\item \textbf{Produits d'exportation}
				\item \textbf{Marges}
				\item \textbf{Accumulation}
				\item \textbf{Total}
			\end{enumerate}
		\end{multicols}
		\label{tab:MACROSAM}%
	\end{table}%
	
	
	\section{Les Ménages}
	
	Un ménage représentatif $h$ perçoit un revenu global $\left(\YH_h\right)$ constitué, d'une part, de la rémunération des facteurs de production : ils captent respectivement une fraction $\psi_h$ de la masse salariale totale et une fraction $\lambda_h$ de la rémunération totale du capital à l'échelle de l'économie. D'autre part, ce revenu est complété par des transferts institutionnels, notamment les pensions versées par l'Administration centrale $\left(\TRA_{h,\gvt}\right)$, les transferts en provenance des autres Administrations Publiques et des ISBLSM $\left(\TRA_{h,\ngo}\right)$, des dividendes versées par les Sociétés $\left(\TRA_{h,\firm}\right)$, ainsi que les transferts courants reçus du Reste du Monde $\left(\TRA_{h,\row}\right)$.
	
	Le revenu disponible ($\YDH_h$) représente la part du revenu total subsistant après déduction des prélèvements obligatoires, notamment les impôts directs ($\TDH_h$). Ce revenu disponible constitue la contrainte budgétaire finale du ménage. Il est d'abord affecté à la consommation finale de biens et services ($\CTH_h$) ainsi qu'aux ressources non fiscales ($\NTXH_h$) versées à l'Administration centrale et aux autres administrations publiques ou institutions assimilées. Ces versements non fiscaux regroupent les redevances et frais administratifs perçus par les administrations décentralisées (Autres APU) d'une part, et les cotisations, redevances de services ou contributions privées perçues par les ISBLSM d'autre part. Enfin, une fraction résiduelle $\psi_h$ du revenu disponible est allouée à l'épargne ($\SH_h$).
	
	L'équilibre ressources-emplois du ménage s'écrit	
	\begin{gather}
		\centerline{\textcolor{Fuchsia}{\textbf{Ressources}:}} \notag\\
		\phi_h \left(e\cdot W^R\cdot\LD^R + \sum_{j}^{}W\cdot \LD_j\right) + \lambda_h\left(e\cdot R^R\cdot\KD^R + \sum_{\trj}^{} R_{\trj}\cdot\KD_{\trj}\right) \notag\\
		+\TR_{h,\firm} + \TR_{h,\gvt} + \TR_{h,\ngo} + e\cdot\TR_{h,\row}\notag\\
		\centerline{\textcolor{Fuchsia}{\textbf{Emplois}:}} \notag\\
		\TDH_h + \NTXH_{h}  + \CTH_h + \SH_h. \notag
	\end{gather}
	
	\section{Les Administrations publiques et les Institutions Sans But Lucratif au Service des Ménages (ISBLSM)}
	\subsection{L'Administration centrale}
	
	L'Administration centrale collecte les impôts indirects sur les produits marchands ($\TI_{tr}$) ainsi que les taxes indirectes additionnelles collectées par les sociétés ($\TIF$). Elle perçoit également les droits de douane ($\TIM_{tr}$) et les prélèvements sur le commerce international versés par les sociétés ($\TIMF$). En matière de fiscalité directe, elle prélève l'impôt sur le revenu des ménages ($\TDH_h$) et celui des sociétés ($\TDF$). Ses recettes non fiscales comprennent celles versées par les ménages ($\NTXHH_{\gvt}$) et les sociétés ($\NTXF$). Son budget est complété par des aides budgétaires ($\TRA_{\gvt,\row}$) provenant de transferts versés par le Reste du monde.
	
	Les dépenses de l'Adminstration centrale comprennent les transferts vers les autres agents, notamment les pensions aux ménages $\left(\TRA_{h,\gvt}\right)$, les transferts privés vers les sociétés $\left(\TRA_{\firm,\gvt}\right)$ qui incluent les intérêts de la dette publique intérieure, les transferts vers les autres APU et les ISBLSM $\left(\TRA_{\ngo,\gvt}\right)$ ainsi que les transferts vers le Reste du monde $\left(\TRA_{\row,\gvt}\right)$ qui comprennent les intérêts de la dette publique extérieure . Les dépenses courantes $G_{\pub}$ servent à produire les services publics. Elles se décomposent en rémunération des salariés de l'État $W\cdot\LD_{\pub}$ ainsi que l'achat de biens et services au titre de la consommation intermédiaire $\left(\PCI_{\pub}\cdot\CI_{\pub}\right)$. L'épargne publique ($\SG$) mesure le solde courant du budget de l'État.
	
	L'équilibre ressources-emplois de l'Administration centrale s'écrit		
	\begin{gather}
		\centerline{\textcolor{Fuchsia}{\textbf{Ressources}:}} \notag\\
		\sum_{\tr}\TI_{\tr} + \TIF + \sum_{\tr}\TIM_{tr} + \TIMF + \TDF +  \sum_{h} \TDH_h\notag\\
		+ \NTXF + \NTXHH_{\gvt} + e\cdot \TR_{\gvt,\row}\notag\\
		\centerline{\textcolor{Fuchsia}{\textbf{Emplois}:}} \notag\\
		\CG_{\pub} + \TR_{\firm,\gvt}  + \sum_{h}\TR_{h,\gvt} + \TR_{\ngo,\gvt} + e\cdot\TR_{\row,\gvt} + \SG.\notag
	\end{gather}
	
	En ce qui concerne la fiscalité directe, nous utilisons les taux effectifs d'imposition appliqués aux différents agents institutionnels. Ainsi, les impôts sur le revenu des ménages et des sociétés sont calculés directement sur la base des revenus perçus. Par ailleurs, nous supposons que les autres taxes à la charge des entreprises sont également indexées sur leurs revenus (\ref{it:TIF} et \ref{it:TIM}). La fiscalité indirecte est modélisée par l'application de taux spécifiques selon la nature et l'origine des flux. Le modèle intègre, d'une part, des taux de droits de douane et de taxes à l'importation ($\tm_{\tr}$) appliqués sur la valeur internationale des produits. D'autre part, il intègre des taux de taxes sur les transactions intérieures ($\tx_{\tr}$) qui s'appliquent non seulement à la production domestique, mais aussi à la valeur des produits importés après acquittement des droits de douane. Enfin, les recettes non fiscales sont modélisées comme des transferts effectués par les sociétés et les ménages dont le montant est calibré comme une proportion de leur revenu disponible et du revenu total (\ref{it:TRAGVTh} et \ref{it:TRAGVTf}, respectivement).
	
	\subsection{Les autres administrations publiques et les ISBLSM}
	Le revenu des autres administrations publiques et des ISBLSM provient principalement des transferts versés par l'Administration centrale ($\TR_{\ngo,\gvt}$), des transferts versés par les sociétés ($\TR_{\ngo,\firm}$), des transferts courants reçus du Reste du Monde ($\TR_{\ngo,\row}$), ainsi que des ressources non fiscales et contributions privées versées par les ménages ($\NTXHH_{\ngo}$).
	
	Ces revenus sont ensuite redistribués sous forme de transferts vers chaque ménage $h$ ($\TRA_{h,\ngo}$). Le solde des revenus est affecté aux dépenses de consommation finale ($G_{\othntr}$), lesquelles se décomposent en rémunération des salariés ($W\cdot\LD_{\othntr}$) et en achats de biens et services ($\PCI_{\othntr}\cdot\CI_{\othntr}$) pour la production des autres biens publics et des autres services non marchands. L'épargne ($\SNGO$) est alors déterminée de manière résiduelle, comme le solde du compte après déduction de l'ensemble des transferts versés et des dépenses de consommation courantes.
	
	L'équilibre ressources-emplois des autres APU et des ISBLSM s'écrit
	\begin{gather}
		\centerline{\textcolor{Fuchsia}{\textbf{Ressources}:}} \notag\\
		\TR_{\ngo,\gvt}+ e\cdot\TR_{\ngo,\row} +\TR_{\ngo,\firm} + \NTXHH_{\ngo} \notag\\
		\centerline{\textcolor{Fuchsia}{\textbf{Emplois}:}} \notag\\
		\sum_{h}\TR_{h,\ngo} + \CG_{\othntr} + \SNGO.\notag
	\end{gather}	
	
	Nous émettons l'hypothèse que la part des dépenses courantes des autres APU et des ISBLSM dans les dépenses courantes publiques totales ($\CG_\pub + \CG_\othntr$) est une fraction $\mu^{\CG}$. Il en résulte ainsi l'équation (\ref{it:Gothntr}). Par ailleurs, en l'absence de données précises, nous supposons que le solde des autres APU et des ISBLSM est maintenu à zéro. En annulant le solde de ces agents secondaires, nous évitons d'introduire une épargne ou un besoin de financement artificiel dans le circuit économique. Cette hypothèse de simplification standard permet ainsi de neutraliser les biais théoriques sur le bouclage macroéconomique et de garantir que les équilibres globaux restent exclusivement pilotés par les comportements des agents majeurs.
	
	\section{Les Sociétés}
	Les Sociétés agissent comme un agent intermédiaire dont le revenu provient principalement de la part de la valeur ajoutée non distribuée aux autres agents. Le revenu total des sociétés $\left(\YF\right)$ est ainsi constitué par la part de l'Excédent Brut d'Exploitation (EBE) qui n'est versée ni aux ménages, ni au Reste du Monde ($\YFK$). À cette base s'ajoutent divers flux de transferts institutionnels ($\YFTR$). La firme reçoit des transferts versés par l'État qui comprennent les intérêts sur la dette publique $\left(\TR_{\firm,\gvt}\right)$. Enfin, les sociétés bénéficient de transferts courants nets en provenance du Reste du Monde $\left(\TR_{\firm,\row}\right)$.
	
	En contrepartie, la firme s'acquitte de diverses obligations fiscales et parafiscales. Elle assure notamment le paiement de l'impôt sur les sociétés ($\TDF$), ainsi que le versement à l'Administration centrale de redevances et autres ressources non fiscales ($\NTXF$). Par ailleurs, bien que collectées pour le compte de l'État, la firme supporte administrativement les autres taxes sur la production ($\TIF$) ainsi que les autres taxes sur le commerce extérieur ($\TIMF$) qui ne sont pas directement appliquées sur les produits. Les sociétés redistribuent également une partie de leur revenu disponible sous forme de transferts vers les ménages $\left(\TRA_{h,\firm}\right)$ et vers les autres administrations publiques et ISBLSM $\left(\TRA_{\ngo,\firm}\right)$. L'épargne des sociétés ($\SF$) correspond au solde du compte de revenu après déduction de l'ensemble des charges fiscales, des ressources non fiscales et des transferts versés, représentant ainsi la capacité d'autofinancement des firmes pour leurs investissements futurs.
	
	L'équilibre ressources-emplois des Sociétés s'écrit:
	\begin{gather}
		\centerline{\textcolor{Fuchsia}{\textbf{Ressources}:}} \notag\\
		\left(1 - \lambda^R-\sum_{h}\lambda_h\right)\cdot\left(e\cdot R^R\cdot\KD^R + \sum_{\trj}^{} R_{\trj}\cdot\KD_{\trj}\right)+ \TR_{\firm,\gvt} + e\cdot\TRA_{\firm,\row} \notag\\
		\centerline{\textcolor{Fuchsia}{\textbf{Emplois}:}} \notag\\
		\TDF + \TIF + \TIMF + \NTXF + \TRA_{\ngo,\firm} + \sum_{h} \TRA_{h,\firm} + \SF.\notag
	\end{gather}
	
	
	
	\section{Le Reste du monde}
	Le compte du Reste du Monde retrace l'ensemble des flux réels et financiers entre l'économie nationale et l'extérieur. Les ressources de cet agent correspondent aux paiements effectués par l'économie domestique vers l'étranger. Elles comprennent la rémunération des travailleurs non-résidents $(\YROWL)$, calculée comme la part de la masse salariale totale non allouée aux travailleurs résidents. Elles incluent également les dividendes versés aux non-résidents $(\YROWK)$, définis comme une fraction $\lambda^R$ du revenu total du capital de l'économie, ainsi que les intérêts sur la dette publique extérieure $\TRA_{\row,\gvt}$. À ces flux s'ajoutent les recettes issues des importations, évaluées aux prix mondiaux et converties par le taux de change $\left(\PWM_{\tr} \cdot \IM_{\tr}\right)$.
	
	En contrepartie, le Reste du Monde effectue divers versements au profit de l'économie nationale. Ces emplois incluent la compensation des employés nationaux travaillant à l'étranger $\left(W^R \cdot \LD^R\right)$ ainsi que les revenus des investissements $\left(R^R \cdot \KD^R\right)$. Le Reste du Monde assure également des transferts courants vers les autres agents institutionnels : les ménages $\left(\TRA_{h,\row}\right)$, les firmes $\left(\TRA_{\firm,\row}\right)$, l'Administration centrale $\left(\TRA_{\gvt,\row}\right)$ et les autres APU et les ISBLSM $\left(\TRA_{\ngo,\row}\right)$. Enfin, cet agent finance l'acquisition des exportations de biens marchands, dont la valeur est déterminée par le prix mondial et le taux de change $\left(\PWE_{\tr} \cdot \EX_{\tr}\right)$. Enfin, l'épargne du Reste du Monde ($\SROW$) est déterminée comme le solde entre ses ressources et ses emplois. Cette épargne correspond au déficit de la balance courante.
	
	Nous adoptons l'hypothèse de petite économie ouverte, en supposant que la taille de l'économie de Madagascar n'est pas suffisamment significative pour influencer les prix mondiaux. En conséquence, les prix des biens exportés $\left(\PWE_{\tr}\right)$ et importés $\left(\PWM_{\tr}\right)$ sont considérés comme exogènes et déterminés sur le marché international.
	
	L'équilibre ressources-emplois du Reste du monde s'écrit (toutes les variables étant exprimées en monnaies étrangères):
	\begin{gather}
		\centerline{\textcolor{Fuchsia}{\textbf{Ressources}:}} \notag\\
		\sum_{\tr}^{} \PWM_{\tr}\cdot\IM_{tr} + \left(1 -\sum_{h}\phi_h\right) \cdot\left(e\cdot W^R\cdot\LD^R + \sum_{j}^{} W\cdot\LD_{j}\right) \notag\\
		+ \lambda^R \cdot\left(R^R\cdot\KD^R + \sum_{\tr}^{} R_{\tr}\cdot\KD_{\tr}\right) + \TR_{\row,\gvt}\notag\\
		\centerline{\textcolor{Fuchsia}{\textbf{Emplois}:}} \notag\\
		\sum_{\tr}^{} \PWE_{\tr}\cdot\EX_{\tr}  +  W^R\cdot\LD^R +  R^R\cdot\KD^R\notag\\
		+ \sum_{h}\TR_{h,\row} + \TR_{\firm,\row} + \TR_{\gvt,\row} + \TR_{\ngo,\row} + \SROW. \notag
	\end{gather}
	
	
	
	\section{Les autres composantes de la demande}
	\subsection{Les dépenses d'investissement}
	Les dépenses d'investissement sont financées par l'épargne totale disponible dans l'économie (\ref{it:IT}). La dépense d'investissement par produit est modélisée comme une part fixe de l'investissement total hors variations de stocks (\ref{it:INV}). Dans ce cadre, le volume total de l'investissement dans l'économie est déterminé de manière endogène par le montant de l'épargne totale disponible. Cette hypothèse reflète fidèlement la situation de Madagascar, où les autorités procèdent régulièrement à des ``coupes'' dans les dépenses d'investissement lorsque l'Administration centrale doit ajuster son budget ou faire face à une contrainte de financement.
	
	\subsection{Modélisation de la demande et de l'équilibre du marché des produits pétroliers}
	
	L'offre disponible de produits pétroliers sur le marché intérieur est entièrement assurée par les importations. Nous posons l'hypothèse selon laquelle l'intégralité du volume importé est immédiatement disponible à la vente $\left(\IM_{\petr} = Q_{\petr}\right)$. Une première fraction $\left(\mu^D\right)$ de cette offre est destinée à satisfaire la demande locale (\ref{it:muDQ}), constituée de la consommation finale des ménages, de la consommation intermédiaire des branches productives et de la demande finale en produits pétroliers pour fins d'investissement. Une seconde fraction $\left(\mu^E\right)$ est allouée à la réexportation (\ref{it:EXpetr}). Enfin, tout écart entre l'offre et les utilisations est résorbé par la variation de stock $\left(\VSTK_{\petr}\right)$, qui joue ainsi le rôle de variable d'ajustement assurant l'équilibre matériel sur le marché pétrolier (\ref{it:QDpetr}).
	
	\subsection{Les marges de commercialisation et de transport}
	Les marges de commercialisation et de transport correspondent aux volumes de services de logistique et de distribution indispensables pour acheminer les biens depuis leur lieu de production ou le cordon douanier jusqu'au consommateur final. Les marges sont ici modélisées comme des parts fixes appliqués aux quantités physiques. Ainsi, des taux de marge sont appliqués au volume de la production nationale et des importations afin de déterminer les quantités de services de marge nécessaires à la distribution des marchandises auprès du consommateur final (\ref{it:MGCOM} et \ref{it:MGTRP}). Ces besoins en services interviennent directement dans la formation des prix en créant un écart entre le prix au producteur et le prix du bien composite. On distingue d'abord les prix des produits domestiques incluant les taxes indirectes ($\PPD_{tr}$) et les prix des produits importés incluant les droits de douane et les taxes indirectes ($\PPM_{tr}$), tous deux définis hors marges (\ref{it:PPD} et \ref{it:PPM}). On définit ensuite les prix des produits domestiques et importés toutes taxes et marges comprises (valables pour les produits biens $\bns$), notés respectivement $\PD_{\tr}$ et $\PM_{\tr}$ (\ref{it:PDbns}, \ref{it:PDser}, \ref{it:PMbns} et \ref{it:PMser}). Dans cette structure, le coût unitaire de la marge dépend du prix du service  ($\PC_{\ser}$) multiplié par les taux de marge en volume. Enfin, la somme de ces quantités de services de marges constitue une composante intégrante de la demande totale adressée au secteur des services (\ref{it:walras}).
	
	
	
	\section{Les fonctions de comportement}
	\subsection{Les ménages}
	Le comportement du ménage $h$ est modélisé par une fonction d'utilité de type Cobb-Douglas. Le ménage cherche à maximiser son niveau de satisfaction global en répartissant son budget entre les différents produits disponibles sur le marché. Le programme d'optimisation s'énonce comme suit:
	\begin{equation}
		\max U_h = \prod_{\tr}^{} \CH_{tr,h}^{\gamma^{\CH}_{\tr,h}},
	\end{equation}
	sous la contrainte budgétaire
	\begin{equation}
		\sum_{\tr} \PC_{\tr}\cdot\CH_{\tr,h} = \CTH_h,
	\end{equation}
	où $\CH_{\tr,h}$ est le volume de la consommation du ménage $h$ pour le produit $\tr$, $\PC_{\tr}$ est le prix à la consommation du produit tr (incluant les taxes et les marges), et $\gamma^{\CH}_{\tr,h}$ est la représente l'élasticité de l'utilité par rapport à la consommation du bien  $\tr$, avec $0\leq \gamma^{\CH}_{\tr,h}\leq 1 $ et $\sum_{\tr}\gamma^{\CH}_{\tr,h}=1$.
	
	L'expression de la demande optimale pour le bien $\tr$  donnée par l'équation \ref{it:CHtrh} stipule que la dépense d'investissement pour le bien $\tr$ est une fraction constante $\gamma^{\CH}_{\tr,h}$ du budget total $\CTH_h$.
	\subsection{Les producteurs}
	La structure de production de la branche $j$ est organisée selon une hiérarchie à plusieurs niveaux (cf.~\textsc{Figure}~\ref{fig:xsvaci}). Chaque secteur de production combine le travail et le capital avec les inputs intermédiaires pour produire un output qui peut être vendu sur le marché domestique exporté.
	%		\incgraph[a3paper, landscape,	overlay page number at bottom][width=1.18\paperheight]{trade}
	
	
	\begin{figure}[!htbp]
		\centering
		\input{Structure_production}
		\caption{Structure de production emboîtée du modèle}
		\label{fig:xsvaci}
	\end{figure}
	
	
	
	Au niveau supérieur, la production totale $\left(\XS_{j}\right)$ est modélisée par une technologie de type Leontief. Cette formulation suppose une complémentarité stricte entre la valeur ajoutée $\left(\VA_{j}\right)$ et l'agrégat des consommations intermédiaires $\left(\CI_{j}\right)$: 
	
	\begin{equation*}
		\XS_{j} = \min\left\{\frac{\VA_{j}}{v_{j}},\frac{\CI_{j}}{\io_{j}}\right\}, 
	\end{equation*}	
	où $v_{j}$ et $\io_{j}$ représentent respectivement les coefficients techniques de valeur ajoutée et de consommation intermédiaire totale par unité produite par la branche $j$.
	
	Au niveau inférieur, la détermination de la valeur ajoutée $\left(\VA_{j}\right)$ diffère selon la nature de la branche. Pour les {branches marchandes} $\trj$, la valeur ajoutée résulte de la combinaison du capital $\left(\KD_{\trj}\right)$ et du travail $\left(\LD_{\trj}\right)$ via une fonction de type Cobb-Douglas. Le producteur y minimise son coût total $C_{\VA}$ :
	\begin{equation*}
		\text{Min } C_{\VA} = R_{\trj} \cdot \KD_{\trj} + W \cdot \LD_{\trj}
	\end{equation*}
	sous la contrainte de production de la valeur ajoutée :
	\begin{equation*}
		\VA_{\trj} = A_{\trj} \cdot \KD_{\trj}^{\alpha_{\trj}} \cdot \LD_{r}^{1-\alpha_{\trj}}.
	\end{equation*}
	Dans cette expression, $R_{\trj}$ désigne le coût du capital, $W$ le taux de salaire moyen, et $\alpha_{\trj}$ l'élasticité de la valeur ajoutée par rapport au capital pour la branche d'activité marchande $\trj$. Le système de solutions issu de ce problème d'optimisation est défini par les équations \ref{it:LD} et \ref{it:KD}.
	
	Pour les branches d'activités non marchandes $\ntr$, la valeur ajoutée est supposée dépendre exclusivement du facteur travail:
	\begin{equation*}
		\VA_{\ntr} =\LD_{\ntr}.
	\end{equation*}
	
	
	Parallèlement, l'agrégat de consommation intermédiaire ($\CI_{j}$) est également régi par une technologie de type Leontief. La demande $\DI_{\tr,j}$ pour chaque produit spécifique $\tr$ utilisé comme intrant par la branche $j$ est strictement proportionnelle au volume total des consommations intermédiaires de ladite branche :
	\begin{equation*}
		\DI_{\tr,j} = \aij_{\tr,j}\cdot \CI_j,
	\end{equation*}
	où $\aij_{\tr,j}$ représente les coefficients de la matrice des entrées intermédiaires, indiquant la quantité du produit $\tr$ nécessaire à la production d'une unité de l'agrégat de consommation intermédiaire de la branche $j$. 
	
	\begin{figure}[!htbp]
		\centering
		\input{commerce_int}
	\end{figure}
	
	%
	Par ailleurs, la stratégie de production de la branche d'activité marchande $\trj$ repose aussi sur une décomposition hiérarchique (cf.~\textsc{Figure}~\ref{fig:relRDM}). Le producteur arbitre sa production totale $\left(\XS_{\trj}\right)$ entre un produit d'exportation ``phare'' ($\EX_{\trj}$) et un agrégat de production locale $\left(\DX_{\trj}\right)$, ce dernier étant structurellement constitué d'un mélange de différents produits spécifiques. Dans ce cadre, le producteur cherche à maximiser son profit en fonction des prix relatifs observés sur les marchés extérieur et domestique. Cet arbitrage est modélisé par une fonction à {Élasticité de Transformation Constante (CET)} définie par l'équation de maximisation suivante :
	
	\begin{equation*}
		\max \Pi = P_{\trj} \cdot \XS_{\trj} - \left( \PLD_{\trj} \cdot \DX_{\trj} + \PE_{\trj} \cdot \EX_{\trj} \right),
	\end{equation*}
	sous la contrainte technique de transformation :
	\begin{equation*}
		\XS_{\trj} = A_{\trj}^E \cdot \left[\beta_{\trj}^E\cdot\EX_{\trj}^{\rho_{\trj}^E} + \left(1 - \beta_{\trj}^E\right)\cdot \DX_{\trj}^{\rho_{\trj}^E}\right]^{\frac{1}{^{\rho_{\trj}^E}}},
	\end{equation*}
	où est $A_{\trj}^E$ est le paramètre d'échelle, $\beta_{\trj}^E$ est le paramètre de répartition qui reflète l'inclinaison du secteur à s'orienter vers les marchés extérieurs par rapport au marché domestique, et $\rho_{\trj}^E$ est le paramètre de transformation,
	L'arbitrage effectué conduit à l'expression de l'équation \ref{it:EX}.
	
	Contrairement au produit phare, la production locale ($\DX_{\trj}$) est un agrégat dont la composition est prédéfinie par une technologie de type Leontief. Chaque produit marchand $\tr$ représente une part fixe de ce volume local, ce qui traduit une spécialisation rigide de l'outil productif pour le marché intérieur. L'offre du produit $\tr$ issue de la branche d'activité marchande $\trj$ est alors donnée par la relation
	\begin{equation*}
		\DS_{\tr} = \theta_{\tr, \trj} \cdot \DX_{\trj},
	\end{equation*}
	où $\theta_{\tr, \trj}$ désigne la part fixe du produit $\tr$ au sein du mélange constituant la production locale de la branche $\trj$. En conséquence, le prix de l'agrégat domestique ($\PLD_{\trj}$) est déterminé par la somme pondérée des prix des produits qui le composent :
	\begin{equation*}
		\PLD_{\trj} = \sum_{\tr} \theta_{\tr, \trj} \cdot \PL_{\tr},
	\end{equation*}
	où $\PL_{\tr}$ désigne le prix du produit $\tr$ sur le marché local. 
	
	
	\subsection{Les consommateurs finaux}
	Le comportement des consommateurs repose sur l'hypothèse d'Armington, selon laquelle les produits domestiques et importés sont considérés comme des substituts imparfaits. Dans cette structure, l'utilisateur final ne distingue pas la branche spécifique $\trj$ ayant produit le bien; il perçoit uniquement une offre locale globale pour chaque produit $\tr$, résultant de l'agrégation de la production de toutes les branches marchandes. 
	
	Le consommateur arbitre ainsi entre cette offre domestique totale ($\DD_{\tr}$) et le produit importé ($\IM_{\tr}$). L'offre domestique totale ($\DD_{tr}$) utilisée dans cet arbitrage correspond à la somme des productions locales de toutes les branches marchandes $\trj$ pour le produit $\tr$ :
	\begin{equation*}
		\DD_{\tr} = \sum_{\trj} \DS_{\tr, \trj}.
	\end{equation*}
	Pour un niveau de demande totale $Q_{\tr}$ donné, il cherche à minimiser son coût d'acquisition. Ce comportement est modélisé par une fonction à {Élasticité de Substitution Constante (CES)}, qui agrège les deux sources d'approvisionnement :
	\begin{equation*}
		\min \PC_{\tr} \cdot Q_{\tr} = \PD_{tr} \cdot \DD_{\tr} + \PM_{\tr} \cdot \IM_{\tr},
	\end{equation*}
	sous la contrainte de la fonction d'agrégation d'Armington
	\begin{equation*}
		Q_{\tr} = A_{\tr}^M \cdot \bigg[\beta_{\tr}^M\cdot\IM_{\tr}^{-\rho_{\tr}^M} + \left(1 - \beta_{\tr}^M\right)\cdot\DD_{\tr}^{-\rho_{\tr}^M}\bigg]^{-\frac{1}{^{\rho_{\tr}^M}}},
	\end{equation*}
	où $A_{\tr}^M$ est un paramètre d'échelle, $\beta_{\tr}^M$ est le paramètre de distribution qui indique la préférence pour les importations dans la consommation totale, $\rho_{\tr}^M$ est le paramètre de substitution qui détermine la facilité avec laquelle l'agent peut remplacer le produit local par le produit importé. La résolution des conditions de premier ordre de ce programme d'optimisation permet d'aboutir à la fonction de demande optimale dont l'expression est donnée par l'équation \ref{it:IM}.
	
	
	
	

	
	
	\section{Intégration du bloc financier}    
	Nous proposons d'intégrer un bloc financier au modèle. Le système du bloc réel étant déjà pleinement déterminé de l'équation \eqref{it:VAj} à l'équation \eqref{it:walras}, le bloc financier fonctionne selon une structure séquentielle (\textit{top-down}). Il prend en entrées les variables endogènes issues du bloc réel ainsi que la valeur des variables exogènes prédéterminées. Le cadre comptable assurant la cohérence d'ensemble de cette extension repose sur la \textit{Matrice de Comptabilité Sociale Financière} (MCSF). Afin de garantir la fermeture comptable et la faisabilité empirique du modèle, le \textit{secteur public} désigne exclusivement le compte de l'Administration Centrale, tandis que le \textit{secteur privé} regroupe par agrégation l'ensemble des autres agents institutionnels résidents, à savoir les ménages, les sociétés non financières, les institutions financières et les ISBLSM.
	
	
	\begin{table}[!h]
		\renewcommand{\arraystretch}{1.5} % Scales row height (1.0 is default)
		\centering
		\caption{Matrice de Comptabilité Sociale Financière de Madagascar. Année 2019. En Milliards d'Ariary.}
		\resizebox{\columnwidth}{!}{%
			\begin{tabular}{|l|r|r|r|r|r|r|r|r|}
				\hline
				&\sffamily \textbf{DFIN}  &\sffamily \textbf{FFIN}  & \textbf{FFDI}  &\sffamily \textbf{GFIN}  &\sffamily \textbf{PFIN}  &\sffamily \textbf{CAPGAIN} &\sffamily \textbf{ACAP}  &\sffamily \textbf{TOTAL} \\
				\hline
				\sffamily \textbf{DFIN}  &       &       &       &       & 835.7 &       &       & \textbf{835.7} \\
				\hline
				\sffamily \textbf{FFIN}  & -511.7 &       &       & 234.4 &       & 0.7   & 1653.2 & \textbf{1376.7} \\
				\hline
				\sffamily \textbf{FFDI}  &       & 1350.4 &       &       & 366.7 &       &       & \textbf{1717.1} \\
				\hline
				\sffamily \textbf{GFIN}  & -25.0 & 5.4   &       &       & 66.7  & 2105.9 & 936.1 & \textbf{3089.0} \\
				\hline
				\sffamily \textbf{PFIN}  & 1173.5 & -2085.0 & 1717.1 &       &       & 198.1 & 8987.6 & \textbf{9991.4} \\
				\hline
				\sffamily \textbf{CAPGAIN} & 198.8 & 2105.9 &       &       &       &       &       & \textbf{2304.7} \\
				\hline
				\sffamily \textbf{ACAP}  &       &       &       & 2854.6 & 8722.3 &       &       & \textbf{11576.9} \\
				\hline
				\sffamily \textbf{TOTAL}   & \textbf{835.7} & \textbf{1376.7 } & \textbf{1717.1} & \textbf{3089.0} & \textbf{9991.4} & \textbf{2304.7} & \textbf{11576.9} &  \\
				\hline
		\end{tabular}}%
		\label{tab:FSAMvalues}%
	\end{table}%
	
	
	% Table generated by Excel2LaTeX from sheet 'Feuil1'
	\begin{table}[!h]
		\renewcommand{\arraystretch}{1.5} % Scales row height (1.0 is default)
		\centering
		\caption{Add caption}
		\resizebox{\columnwidth}{!}{%
			\begin{tabular}{|l|c|c|c|c|c|c|c|}
				\hline
				&\sffamily \textbf{DFIN}  &\sffamily \textbf{FFIN}  & \textbf{FFDI}  &\sffamily \textbf{GFIN}  &\sffamily \textbf{PFIN}  &\sffamily \textbf{CAPGAIN} &\sffamily \textbf{ACAP} \\[-1.5ex]
				&\sffamily \textbf{(1)} &\sffamily \textbf{(2)}&\sffamily \textbf{(3)} &\sffamily \textbf{(4)}&\sffamily \textbf{(5) }&\sffamily \textbf{(6)}&\sffamily \textbf{(7)}\\
				\hline
				\sffamily\textbf{(1)}  &       &       &       &       & $\Delta\MD$ &       &       \\
				\hline
				\sffamily\textbf{(2)}  & $\ER\cdot\Delta\RES$ &       &       & $e\cdot\AMORTG$ &       & $\left(e - \ER\right)\cdot\Delta\RES$   & $e\cdot\SROW$  \\
				\hline
				\sffamily\textbf{(3)}  &       & $e\cdot\FDI^{\mathit{NET}}$ &       &       & $e\cdot\FDI^{\mathit{OUT}}$ &       &        \\
				\hline
				\sffamily\textbf{(4)}  & $\Delta\DCG$ & $e\cdot\Delta\NFDG^{\mathit{prg}}$   &       &       & $\Delta\BG$  & $\begin{matrix}e\cdot\KTRG\\[-1.5ex] +\\[-1.5ex] e\cdot\Delta\NFDG^{\mathit{prj}}\end{matrix}$ & $\SG$ \\
				\hline
				\sffamily\textbf{(5)}  & $\Delta\DCP$ & $e\cdot\Delta\NFDP$ & $e\cdot\FDI^{\mathit{IN}}$ &       &       & $
				\begin{matrix} e\cdot\OTHKTR\\[-1.5ex] +\\[-1.5ex] \RES_{\textit{-1}}\cdot\Delta \ER\\[-1.5ex] +\\[-1.5ex] (\ER - e)\cdot\Delta\RES \end{matrix}$ & $\SP
				$  \\
				\hline
				\sffamily\textbf{(6)} & $\RES_{\textit{-1}}\cdot\Delta\ER$ & $\begin{matrix}
					e\cdot\OTHKTR \\[-1.5ex]
					+\\[-1.5ex]
					e\cdot\KTRG\\[-1.5ex]
					+\\[-1.5ex]
					e\cdot\Delta\NFDG^{\mathit{prj}}
				\end{matrix}$ &       &       &       &       &        \\
				\hline
				\sffamily\textbf{(7)}  &       &       &       & $\INVG$ & $\INVP$ &       &        \\
				\hline
		\end{tabular}}%
		\label{tab:FSAMformula}%
	\end{table}%
	
	\subsection{Matrice de Comptabilité Sociale Financière}
	La MCSF utilisée dans ce projet est une version légèrement modifiée de la version agrégée présentée dans les travaux précédents de \cite{Ramaharo2022}, rendue compatible avec la MACROSAM retenue dans ce travail. 
	
	Les différents comptes structurant la matrice financière sont les suivants :
	\begin{itemize}
		\item \textbf{DFIN} : Le compte financier intérieur qui retrace les avoirs et engagements monétaires et bancaires intérieurs.
		\item \textbf{FFIN} : Le compte financier extérieur.
		\item \textbf{FFDI} : Le compte des investissements directs étrangers (IDE).
		\item \textbf{GFIN} et \textbf{PFIN} : Les comptes financiers respectifs du secteur public et du secteur privé.
		\item \textbf{CAPGAIN} : Le compte de gains en capital, compte virtuel destiné à enregistrer les réévaluations d'actifs.
		\item \textbf{ACAP} : Le compte d'accumulation ou d'investissement global.
	\end{itemize}
	
	\subsection{Équilibre financier du secteur public}
	Le besoin de financement net du secteur public ($\BRG$), mesuré par l'écart entre l'investissement et l'épargne publics, est couvert par la combinaison des capitaux reçus de l'extérieur ainsi que des financements extérieurs et intérieurs nets :
	\begin{equation*}
		\INVG - \SG = e \cdot \KTRG + e \cdot \left(\Delta\NFDG^{\mathit{prj}} +  \Delta\NFDG^{\mathit{prg}} - \AMORTG\right) + \Delta\DCG + \Delta\BG
	\end{equation*}
	Dans cette identité, $\INVG$ désigne l'investissement public et $\SG$ l'épargne publique. La variable $\KTRG$ désigne les transferts en capital reçus par le secteur public, tandis que $\Delta\NFDG^{\mathit{prj}}$ et $\Delta\NFDG^{\mathit{prg}}$ mesurent respectivement les tirages sur projets et sur programmes, et $\AMORTG$ représente l'amortissement de la dette extérieure publique. Le financement intérieur se compose du crédit bancaire net à l'État ($\Delta\DCG$) et des émissions nettes de titres publics souscrits par le secteur privé non bancaire ($\Delta\BG$). Le bouclage financier du secteur public s'effectue par le flux de crédit bancaire ($\Delta\DCG$), qui agit comme variable d'ajustement, traduisant le rôle de prêteur en dernier ressort du système bancaire (incluant la Banque Centrale). Par ailleurs, les souscriptions non bancaires sont directement prélevées sur l'épargne privée ($\SP$) selon la règle comportementale :
	\begin{equation*}
		\Delta\BG = \psi^{\BG} \cdot \SP,
	\end{equation*}
	où $\SP = \SF + \SNG + \sum_{h}\SH_h $ et $\psi^{\BG}$ représente la propension marginale du secteur privé à allouer son épargne aux titres publics.   
	

Par ailleurs, nous intégrons un bloc dédié à la dette publique. Le stock total de la dette publique ($\DEBTG$) est déterminé par l'identité comptable suivante :
\begin{equation}
	\DEBTG = \ER \cdot \NFDG + \NFDP,
\end{equation} 
où $\NFDG$ désigne le stock de la dette publique extérieure en devises, tandis que $\NFDP$ designe le stock de la dette publique intérieure. Conformément à la contrainte de soutenabilité extérieure, le secteur public contracte des engagements extérieurs dont l'encours maximal est ancré sur la capacité d'exportation de l'économie. Ainsi, le stock de dette extérieure est modélisé comme une proportion directe de la valeur totale des exportations de biens et services en devises :
\begin{equation*}
	\NFDG = \psi^{\EX} \cdot \sum_{\tr} \PWE_{\tr} \cdot \EX_{\tr},
\end{equation*} 
où $\psi^{\EX}$ représente le ratio d'endettement extérieur rapporté aux exportations, et le terme $\sum_{\tr} \PWE_{\tr} \cdot \EX_{\tr}$ capture la valeur totale des exportations au prix du marché mondial.


 
	
	\subsection{Équilibre du marché monétaire}
	L'équilibre du marché monétaire repose sur l'égalité stricte entre l'offre et la demande de monnaie. La demande de monnaie est formalisée selon une approche d'encaisses de transaction dérivée de la théorie quantitative de la monnaie :
	\begin{equation*}
		\MD = \frac{1}{\vel} \cdot \PIB^{\mathit{MP}},
	\end{equation*}
	où $\vel$ représente la vitesse de circulation de la monnaie et $\PIB^{\mathit{MP}}$ le produit intérieur brut nominal au prix du marché.
	
	Du côté de l'offre, la masse monétaire s'ajuste pour satisfaire cette demande et équivaut à la somme des contreparties monétaires inscrites au bilan consolidé du système bancaire :
	\begin{equation*}
		\MS = \MD = \ER \cdot \RES + \DCG + \DCP
	\end{equation*}
	Dans cette égalité, $\MS$ et $\MD$ désignent respectivement l'offre et la demande de monnaie, et $\RES$ représente le stock d'avoirs extérieurs nets (réserves) en devises. Le crédit intérieur net se décompose entre le stock de créances nettes du système bancaire sur le secteur public ($\DCG$) et le stock de crédit accordé au secteur privé ($\DCP$). 
	En terme de flux, l'identité peut s'écrire
	\begin{equation*}
		\Delta \MD = e\cdot\Delta\RES + \Delta\DCG + \Delta \DCP + \RES_{\textit{-1}}\cdot\Delta\ER + (\ER - e)\cdot\Delta\RES,
	\end{equation*}
	où $e\cdot\Delta\RES$ décrit le flux de transaction, et la quantité $\RES_{\textit{-1}}\cdot\Delta\ER + (\ER - e)\cdot\Delta\RES$ décrit l'ajustement de valeur. Dans cette notation, $\RES_{\textit{-1}}$ correspond au stock de réserves de la période précédant l'année de base.
	Le flux de crédit accordé au secteur privé ($\Delta\DCP$) est directement modélisé en fonction du niveau de l'investissement privé ($\INVP$) :
	\begin{equation*}
		\Delta\DCP = \psi^{\DCP} \cdot \INVP,
	\end{equation*}
	où $\psi^{\DCP}$ représente la part de la formation brute de capital fixe du secteur privé couverte par le crédit bancaire. L'ajustement du marché monétaire s'effectue par la variation des réserves internationales de change ($\Delta\RES$), qui constitue la variable résiduelle du système.
	
	\subsection{Équilibre du marché extérieur et de la balance des paiements}
	L'équilibre du marché extérieur et la dynamique des avoirs extérieurs nets découlent directement de l'identité comptable de la balance des paiements : 
	\begin{equation*}
		\Delta\RES = \CAB + \KTRG + \OTHKTR + \FDI^{\textit{NET}} + \Delta\NFDG^{\mathit{prg}} + \Delta\NFDG^{\mathit{prj}} - \AMORTG + \Delta\NFDP.
	\end{equation*}
	Au sein de cette identité, $\CAB = -\SROW$ désigne la balance courante, tandis que $\OTHKTR$ représente les flux nets de transferts de capital reçus par le secteur privé. Les entrées nettes d'investissements directs étrangers ($\FDI^{\mathit{NET}}$) sont définies par :
	\begin{equation*}
		\FDI^{\mathit{NET}} = \FDI^{\mathit{IN}} - \FDI^{\mathit{OUT}},
	\end{equation*}
	où $\FDI^{\mathit{IN}}$ et $\FDI^{\mathit{OUT}}$ représentent respectivement les IDE entrants et sortants. Par ailleurs, la variable $\Delta\NFDP$ correspond au flux d'endettement extérieur net souscrit par le secteur privé. Dans cette architecture de modélisation, le financement extérieur net du secteur privé ($\Delta\NFDP$) constitue la variable résiduelle d'ajustement de la balance des paiements. Par convention de bouclage, cette variable absorbe l'écart global de financement extérieur et intègre implicitement le poste statistique des erreurs et omissions nettes, garantissant ainsi la stricte cohérence macroéconomique du système.
	
	\subsection{Équilibre résiduel du secteur privé}
	En vertu de la loi de Walras, l'équilibre financier du secteur privé est redondant. Sa capacité ou son besoin de financement net ne constitue pas une équation indépendante, car cette identité se déduit directement de la combinaison des contraintes budgétaires précédentes :
	\begin{equation*}
		\INVP - \SP = \Delta\DCP + e \cdot \left(\OTHKTR + \FDI^{\textit{NET}} + \Delta\NFDP\right) + \RES_{\textit{-1}} \cdot \Delta\ER + (\ER - e) \cdot \Delta\RES - \Delta\MD - \Delta\BG,
	\end{equation*}
	toutes les variables composant cette identité ayant été définies dans les sections précédentes.
	
	
	%    \end{description}

%\begin{equation}
%	\YH_h = \lambda_h\sum_{j}^{} W_j\LD_j + \phi_h\sum_{j}^{} R_j\KD_j + \TG_h + \TR^{ROW}_h + \TR^{NGO}_h
%\end{equation}
%
%%\begin{equation}
%%	\YDH_h = \YH_h - (\TD_h + \RNF_h)
%%\end{equation}
%%
%
%
%\begin{table}[H]
%	\begin{tabular}{lllll}
	%	\toprule
	%	\multicolumn{2}{c}{\textbf{Ressources}} && \multicolumn{2}{c}{\textbf{Emplois}}\\
	%	\midrule
	%Salaire	& $\lambda_h\sum_{j}^{} W\cdot\LD_j$ &  & Impots sur les revenus & $\DTH_h$  \\
	%Revenu du capital	& $\lambda_h\sum_{j}^{} R_j\cdot\KD_j$ &  &  Recettes non fiscales &  $\RNFH_h$\\
	%Transferts recu des ISBLSM & $\TRH^{NG}_h$ & & Dépenses en produits $ \tr$ & $C_{\tr,h}$\\
	%Transferts recu du Reste du monde & $\TRH^{R}_h$ && Epargnes des ménages & $\SH_h$\\
	%\bottomrule
	%%\hline
	%%\textbf{Total des ressources}
	%%	\hline
	%%	&  &  &  &  \\
	%	\hline
	%\end{tabular}
	%\end{table}
	%
	%
	%\section{L'administration centrale}
	%L'administration centrale collecte les impôts indirects sur les produits (TVA, accises), les droits de douane et les autres taxes sur le commerce international. Elle perçoit également les impôts directs sur le revenu des ménages et des sociétés, ainsi que diverses recettes non fiscales. Ces ressources sont ensuite utilisées pour financer la consommation publique, qui se décompose en achats de biens et services (consommations intermédiaires) et en rémunération des salariés de la fonction publique. Ce processus permet la production de services non marchands, l'excédent des recettes sur les dépenses courantes constituant l'épargne publique.
	%
	%\begin{table}[H]
	%	\begin{tabular}{lllll}
		%		\toprule
		%		\multicolumn{2}{c}{\textbf{Ressources}} && \multicolumn{2}{c}{\textbf{Emplois}}\\
		%		\midrule
		%		Taxes directs	      & $\DTH_h + \DTF$         &  & Intérêt sur dettes int         & $\INTD$  \\
		%		Taxe indirect	      & $\TI_{\tr}+\TIF$         &  & Intérêt sur dettes EXT         & $\INTF$  \\
		%		Taxe sur le commerce  & $\TIM_{\tr} + \TIMF_{\tr}$        &  & Transferts vers les autres APU & $\TRNG^G$\\
		%		Recettes non fiscales & $\RNFH_h+ \RNFF$ &  & Dépenses courante          & $\XS_g$\\
		%		Dons courants         & $\TRG^R$      &  & Achats de biens et svc         & $\BSG $\\
		%		                      &               &  & Épargne publique                & $\SG$\\
		%		\bottomrule
		%		%\hline
		%		%\textbf{Total des ressources}
		%		%	\hline
		%		%	&  &  &  &  \\
		%		\hline
		%	\end{tabular}
	%\end{table}
	%
	%Les tax
	%\section{Autres APU et ISBLSM}
	%Les autres APU et ISBLSM receoinvent des trqnsferts due l'administration centrale  et du Reste du monde, verse des transferts aux menages et depenses en produits $i$ et fais des epargnes
	%
	%\begin{table}[H]
	%	\begin{tabular}{lllll}
		%		\toprule
		%		\multicolumn{2}{c}{\textbf{Ressources}} && \multicolumn{2}{c}{\textbf{Emplois}}\\
		%		\midrule
		%		Transferts du Reste du monde	      	& $\TRNG^{R}$         	&  & Transferts vers les menages    & $\TRH^{NG}_h$  \\
		%		Transferts de l'ADM central           	& $\TRNG^{G}$          	&  & Consommation en produits `service' & $C_{\ngo,\ser}$   \\
		%												& 						&  & Epargnes des NGO & $SNGO$     \\
		%														\bottomrule
		%		%\hline
		%		%\textbf{Total des ressources}
		%		%	\hline
		%		%	&  &  &  &  \\
		%		\hline
		%	\end{tabular}
	%\end{table}
	%
	%\section{Les entreprises}
	%Les entreprises recoivent une partie des revenus des facteurs de capital aussi bien domestique que venant du Reste du monde. Ils recoivent aussi une remboursemen du TVA. Il payenet les impots directs sur les revenus, et autres taxes sur le commerc international aisnsi que des recettes non fiscales et fais des epargnes.
	%
	%
	%
	%\begin{table}[H]
	%	\begin{tabular}{lllll}
		%		\toprule
		%		\multicolumn{2}{c}{\textbf{Ressources}} && \multicolumn{2}{c}{\textbf{Emplois}}\\
		%		\midrule
		%		Transferts 	      	& $\TRF^{R}$         	&  & Autres taxes indirects    & $\TIF$  \\
		%		Revenu du capital             	        & $\left(1 - \lambda^R-\sum_{h}\lambda_h\right)\sum_{j}^{} R_j\cdot\KD_j$           	&  & Autres tarifs douaniers & $\TIMF$   \\
		%		& 						&  & Impots sur le revenu & $\DTF$     \\
		%		& 						&  & Recettes non fiscales & $\RNFF$     \\
		%		& 						&  & Impots sur le revenu & $\DTF$     \\
		%		& 						&  & Epargnes & $\SF$     \\
		%		\bottomrule
		%		%\hline
		%		%\textbf{Total des ressources}
		%		%	\hline
		%		%	&  &  &  &  \\
		%		\hline
		%	\end{tabular}
	%\end{table}
	%
	%\section{Le Reste du monde}
	%
	%\textbf{Ressources}] 
	%\begin{itemize}
	%	\item Dividende	      	$\lambda^R\sum_{j}^{} R_j\cdot\KD_j$
	%\end{itemize}
	%
	%\begin{table}[H]
	%	\begin{tabular}{lllll}
		%		\toprule
		%		\multicolumn{2}{c}{\textbf{Ressources}} && \multicolumn{2}{c}{\textbf{Emplois}}\\
		%		\midrule
		%		Dividende	      	& $\lambda^R\sum_{j}^{} R_j\cdot\KD_j$                                               	&  & Transferts vers les menages        & $\TRH^{R}_h$  \\
		%		Compensation des employés	      	& $\left(1 -\sum_{h}\phi_h\right) \sum_{\tr}^{} W\cdot\LD_{j}$         	&  & Transferts vers l'adm. centrale    & $\TRG^{R}$  \\
		%		Interet sur le dette exteerieur          	& $\INTF$          	                                            &  & Transferts vers les APU et ISBLSM  & $\TRNG^{R}$  \\
		%		Importations de biens et services          	& $\IM_{\tr}$          	                                        &  & Transferts vers les entreprises    & $\TRF^{R}$  \\
		%		                                            &                                                               &  & Exportations de biens et services  & $\EX_{\tr}$ \\
		%		                                            &																&  & Balance courante                   & $\CAB$\\
		%		\bottomrule
		%		%\hline
		%		%\textbf{Total des ressources}
		%		%	\hline
		%		%	&  &  &  &  \\
		%		\hline
		%	\end{tabular}
	%\end{table}
	%
	
	
	\section{Spécification du modèle}
	\subsection{Définition des ensembles}
	\begin{itemize}
		\item[$\bullet$] \textbf{Tous les produits}: $i \in \mathbf{I} = \left\{\agr, \food, \othind, \petr, \ser, \pub, \ontr\right\}$
		\begin{description}[labelwidth=\widthof{\bfseries OTHIND}]
			\item[$\agr$]	Agriculture
			\item[$\food$] 	Alimentaire, Boisson, Tabac
			\item[$\othind$] Autres industries
			\item[$\petr$] 	Produits pétroliers
			\item[$\ser$] 	Service marchand
			\item[$\pub$] 	Administration publique
			\item[$\othntr$] 	Autres services non marchands
		\end{description}
		
		\item[$\bullet$] \textbf{Toutes les branches d'activités}:	$j\in \mathbf{J} = \left\{\agr, \food, \othind, \ser, \pub, \ontr\right\}\subset\mathbf{I}$		
		\begin{description}[labelwidth=\widthof{\bfseries OTHIND}]
			\item[$\agr$]	Agriculture
			\item[$\food$] 	Alimentaire, Boisson, Tabac
			\item[$\othind$] Autres industries
			\item[$\ser$] 	Service marchand
			\item[$\pub$] 	Administration publique
			\item[$\othntr$] 	Autres service publiques et services non marchands
		\end{description}
		
		\item[$\bullet$]  \textbf{Les produits marchands}:	$\tr\in\mathbf{TR} = \left\{\agr, \food, \othind, \petr, \ser \right\}\subset\mathbf{I}$
		\begin{description}[labelwidth=\widthof{\bfseries OTHIND}]
			\item[$\agr$]	Agriculture
			\item[$\food$] 	Alimentaire, Boisson, Tabac
			\item[$\othind$] Autres industries
			\item[$\petr$] 	Produits pétroliers
			\item[$\ser$] 	Service marchand
		\end{description}
		
		
		\item[$\bullet$]  \textbf{Les biens marchands}:	$\bns\in\mathbf{TR} = \left\{\agr, \food, \othind, \petr \right\}\subset\mathbf{TR}$
		\begin{description}[labelwidth=\widthof{\bfseries OTHIND}]
			\item[$\agr$]	Agriculture
			\item[$\food$] 	Alimentaire, Boisson, Tabac
			\item[$\othind$] Autres industries
			\item[$\petr$] 	Produits pétroliers
		\end{description}
		
		\item[$\bullet$]  \textbf{Les biens marchands (hors produits pétroliers)}:	$\npbns\in\mathbf{NPBNS} = \left\{\agr, \food, \othind \right\}\subset\mathbf{TR}$
		\begin{description}[labelwidth=\widthof{\bfseries OTHIND}]
			\item[$\agr$]	Agriculture
			\item[$\food$] 	Alimentaire, Boisson, Tabac
			\item[$\othind$] Autres industries
		\end{description}
		
		\item[$\bullet$]  \textbf{Sous-ensemble des produits marchands et des branches d'activités marchandes}:\\	$\trj\in\mathbf{TRJ} = \left\{\agr, \food, \othind, \ser \right\}\subset\mathbf{TR}\subset\mathbf{I}$
		\begin{description}[labelwidth=\widthof{\bfseries OTHIND}]
			\item[$\agr$]	Agriculture
			\item[$\food$] 	Alimentaire, Boisson, Tabac
			\item[$\othind$] Autres industries
			\item[$\ser$] 	Service marchand
		\end{description}
		
		\item[$\bullet$]  \textbf{Les Services non-marchands}:	$\ntr\in\mathbf{NTR} = \left\{\pub, \ontr\right\}\subset\mathbf{J}$
		\begin{description}[labelwidth=\widthof{\bfseries OTHIND}]
			\item[$\pub$] Administration publique
			\item[$\othntr$] Autres services non marchands
		\end{description}
		
		\item[$\bullet$]  \textbf{Les Agents économiques}:	$\ag,\agj\in\mathbf{AG} = \left\{\hur, \hru, \gvt, \ngo, \firm, \row\right\}$
		\begin{description}[labelwidth=\widthof{\bfseries OTHIND}]
			\item[$\hur$] Ménages urbains
			\item[$\hru$] Ménages ruraux
			\item[$\gvt$] Administration centrale
			\item[$\ngo$] Autres Administrations Publiques et ISBLSM
			\item[$\firm$] Sociétés
			\item[$\row$] Reste du monde
		\end{description}
		
		\item[$\bullet$]  \textbf{Les Agents domestiques}:	$\agd\in\mathbf{AGD} = \left\{\hur, \hru, \gvt, \ngo, \firm\right\}\subset\mathbf{AG}$
		\begin{description}[labelwidth=\widthof{\bfseries OTHIND}]
			\item[$\hur$] Ménages urbains
			\item[$\hru$] Ménages ruraux
			\item[$\gvt$] Administration centrale
			\item[$\ngo$] Autres Administrations Publiques et ISBLSM
			\item[$\firm$] Sociétés
		\end{description}
		
		\item[$\bullet$]  \textbf{Les Agents publics}:	$\agpub\in\mathbf{AGPUB} = \left\{\gvt, \ngo, \right\}\subset\mathbf{AGD}$
		\begin{description}[labelwidth=\widthof{\bfseries OTHIND}]
			\item[$\gvt$] Administration centrale
			\item[$\ngo$] Autres Administrations Publiques et ISBLSM
			
		\end{description}
		
		\item[$\bullet$]  \textbf{Les Agents non gouvernementaux}:	$\agngvt\in\mathbf{AGNGVT} = \left\{\hur, \hru, \ngo, \firm, \row\right\}\subset\mathbf{AG}$
		\begin{description}[labelwidth=\widthof{\bfseries OTHIND}]
			\item[$\hur$] Ménages urbains
			\item[$\hru$] Ménages ruraux
			\item[$\ngo$] Autres Administrations Publiques et ISBLSM
			\item[$\firm$] Sociétés
			\item[$\row$] Reste du monde
		\end{description}
		
		\item[$\bullet$] \textbf{Les Ménages}: $h\in \mathbf{H} = \left\{\hur, \hru\right\}\subset\mathbf{AG}$
		\begin{description}[labelwidth=\widthof{\bfseries OTHIND}]
			\item[$\hur$] Ménages urbains
			\item[$\hru$] Ménages ruraux
		\end{description}
		
		
	\end{itemize}
	
	
	
	
	
	
	
	
	
	
	
	\subsection{Liste des équations}
	
	\redparagraph{Bloc production}\hfill \textcolor{Fuchsia}{\textbf{78}}
	\begin{enumerate}[itemsep=2pt,label=\textbf{eq\arabic*}.,ref=\textbf{eq\arabic*}]
		\item $\VA_j = v_j\cdot\XS_j$ \label{it:VAj} \hfill 6
		
		\item $\CI_j = \io_j\cdot\XS_j $ \hfill 6
		
		\item $\VA_{\trj} = A_{\trj}\cdot\LD_{\trj}^{\alpha_{\trj}}\cdot\KD_{\trj}^{1-\alpha_{\trj}}$ \hfill 4
		
		\item $\VA_{\ntr} = \LD_{\ntr} $ \hfill 2
		
		\item $\PVA_{\ntr}\cdot\VA_{\ntr} = W\cdot\LD_{\ntr} $ \hfill 2
		
		\item $W\cdot\LD_{\trj} = \alpha_{\trj}\cdot\PVA_{\trj}\cdot\VA_{\trj}$  \label{it:LD} \hfill 4
		
		\item $R_{\trj}\cdot\KD_{\trj} = (1-\alpha_{\trj})\cdot\PVA_{\trj}\cdot\VA_{\trj}$ \label{it:KD}  \hfill 4
		
		\item $\DI_{\tr,j} = \aij_{\tr,j}\cdot \CI_j$ \hfill 30
		
		\item $\DS_{\tr,\trj} = \theta_{\tr,\trj}\cdot \DX_{\trj}$ \hfill 20
	\end{enumerate}
	
	
	
	\redparagraph{Bloc ménage}\hfill \textcolor{Fuchsia}{\textbf{16}}
	\begin{enumerate}[itemsep=2pt,label=\textbf{eq\arabic*}.,resume,ref=\textbf{eq\arabic*}]
		
		\item $\YH_h  =  \YHL_h + \YHK_h + \YHTR_h $ \hfill 2
		
		\item $\YHL_h = \phi_h \cdot\left(e\cdot W^R\cdot\LD^R + \sum_{j}^{}W\cdot \LD_j\right)$ \hfill 2
		
		\item $\YHK_h = \lambda_h\cdot\left(e\cdot R^R\cdot\KD^R + \sum_{\trj}^{} R_{\trj}\cdot\KD_{\trj}\right) $ \hfill 2
		
		\item $\YHTR_h = \TRA_{h,\firm} + \TRA_{h,\gvt} + \TRA_{h,\ngo} + e\cdot\TRA_{h,\row}$ \hfill 2
		
		\item $\YDH_h = \YH_h - \TDH_h$ \hfill 2
		
		\item $\SH_h = \psi_h\cdot\YDH_h$ \hfill 2
		
		\item $\NTXH_{h} = \ntyh_{h} \cdot\YDH_h$ \label{it:TRAGVTh} \hfill 2
		
		\item $\CTH_h = \YDH_h - \SH_h - \NTXH_{h}$  \hfill 2
		
	\end{enumerate}
	
	
	\redparagraph{Bloc Sociétés}\hfill \textcolor{Fuchsia}{\textbf{5}}
	\begin{enumerate}[itemsep=2pt,label=\textbf{eq\arabic*}.,resume, ref=\textbf{eq\arabic*}]
		
		\item $\YF = \YFK + \YFTR$ \hfill 1
		
		\item $\YFK = \left(1 - \lambda^R-\sum_{h}\lambda_h\right)\cdot\left(e\cdot R^R\cdot\KD^R + \sum_{\trj}^{} R_{\trj}\cdot\KD_{\trj}\right)$ \hfill 1
		
		\item $\YFTR = \TRA_{\firm,\gvt}  + e\cdot\TRA_{\firm,\row} $ \hfill 1
		
		\item $\YDF = \YF  - \TDF - \TIMF -\TIF  $ \hfill 1
		
		\item $\SF = \YDF -  \NTXF  - \TRA_{\ngo,\firm} - \sum_{h}\TRA_{h,\firm}$ \hfill 1
		
	\end{enumerate}
	
	
	\redparagraph{Bloc Administration centrale}\hfill \textcolor{Fuchsia}{\textbf{23}}
	\begin{enumerate}[itemsep=2pt,label=\textbf{eq\arabic*}.,resume]
		
		\item $\YG = \YGTX + \YGNTX + \YGTR$   \hfill 1 
		
		\item $\YGTX = \sum_{\tr}\TI_{\tr} + \TIF + \sum_{\tr}\TIM_{\tr} + \TIMF + \TDF + \sum_{h} \TDH_h $ \hfill 1
		
		\item $\YGNTX = \NTXF + \NTXHH_{\gvt}$   \hfill 1 
		
		\item $\YGTR = e\cdot\TRA_{\gvt, \row} $ \hfill 1
		
		\item $\SG = \YG - \CG_{\pub}  - \sum_{\agngvt\neq\row} \TR_{\agngvt,\gvt} - e\cdot\TR_{\row,\gvt}$ \hfill 1
		
		\item $\TDF = \tyf\cdot\YF$ \hfill 1
		
		\item $\TDH_h  = \tyh_h\cdot\YH_h$ \hfill 2
		
		\item $\begin{aligned}[t]
			\TI_{\bns} &= \tx_{\bns}\cdot \Bigg\{\Big[\PL_{\bns}\ +\left(\tmgc_{\bns} + \tmgt_{\bns}\right)\cdot\PC_{\ser}\Big]\cdot\DD_{\bns}\\
			& + \Big[(1 + \tm_{\bns})\cdot e \cdot \PWM_{\bns} +\left(\tmgc_{\bns} + \tmgt_{\bns}\right)\cdot\PC_{\ser} \Big]\cdot\IM_{\bns} \Bigg\} 
		\end{aligned}$ \hfill 4
		
		\item $\TI_{\ser} = \tx_{\ser}\cdot \Big[\PL_{\ser}\cdot\DD_{\ser} + (1 + \tm_{\ser})\cdot e \cdot \PWM_{\ser}\cdot\IM_{\ser} \Big]$ \hfill 1
		
		\item $\TIF = \syfi\cdot\YF$ \label{it:TIF}  \hfill 1
		
		\item $\TIM_{\tr} = \tm_{\tr} \cdot e \cdot \PWM_{\tr}\cdot\IM_{\tr}$ \hfill 5
		
		\item $\TIMF = \syfm\cdot\YF$  \label{it:TIM}\hfill 1		
		\item $\NTXF = \ntyf \cdot\YDF$ \label{it:TRAGVTf} \hfill 2
		
		
		
		\item  $\NTXHH_{\gvt} = \sntx_{\gvt} \cdot \sum_{h}\NTXH_h$  \hfill 1
	\end{enumerate}
	
	
	\redparagraph{Bloc Autres APU et ISBLSM }\hfill \textcolor{Fuchsia}{\textbf{4}}
	\begin{enumerate}[itemsep=2pt,label=\textbf{eq\arabic*}.,resume,  ref=\textbf{eq\arabic*}]
		
		\item $\YNGO = \TR_{\ngo,\firm} + \TR_{\ngo,\gvt} + e\cdot\TR_{\ngo,\row} + \NTXHH_{\ngo} $ \hfill 1
		
		\item  $\NTXHH_{\ngo} = \sntx_{\ngo} \cdot \sum_{h}\NTXH_h$ \hfill 1
		
		\item $\SNGO = \YNGO - \sum_{h}\TR_{h,\ngo}  - \CG_{\othntr}$ \hfill 1	
		
		\item $\CG_{\othntr} = \left(\dfrac{m^{\CG}}{1-m^{\CG}}\right)\cdot \CG_{\pub}$	\label{it:Gothntr}	\hfill 1		
		
	\end{enumerate}
	
	\redparagraph{Bloc Reste du Monde}\hfill \textcolor{Fuchsia}{\textbf{5}}
	\begin{enumerate}[itemsep=2pt,label=\textbf{eq\arabic*}.,resume,  ref=\textbf{eq\arabic*}]
		
		\item $\YROW =  \YROWL + \YROWK +\sum_{\tr}^{} \PWM_{\tr}\cdot\IM_{tr} + \TR_{\row,\gvt}$ \hfill 1
		
		\item $e\cdot\YROWL = \left(1 -\sum_{h}\phi_h\right)\cdot \left(e\cdot  W^R\cdot\LD^R + \sum_{j}^{} W\cdot\LD_{j}\right) $\hfill 1
		
		\item $e\cdot\YROWK = \lambda^R\cdot \left(e\cdot R^R\cdot\KD^R + \sum_{\trj}^{} R_{\trj}\cdot\KD_{\trj}\right) $\hfill 1
		
		\item $\SROW = \YROW - \sum_{\tr}^{} \PWE_{\tr}\cdot\EX_{\tr}  -  W^R\cdot\LD^R  -  R^R\cdot\KD^R  - \sum_{\agd}\TR_{\agd,\row}$\hfill 1	
		
		\item $\CAB = -\SROW$ \hfill 1
		
	\end{enumerate}
	
	\redparagraph{Bloc Transferts entre agents}\hfill \textcolor{Fuchsia}{\textbf{15}}
	\begin{enumerate}[itemsep=2pt,label=\textbf{eq\arabic*}.,resume,  ref=\textbf{eq\arabic*}]
		
		\item $\TR_{\agd,\row} = \PIXCON^\eta\cdot\TR_{\agd,\row}^{O}$ \hfill 5
		
		\item $\TR_{\agngvt,\gvt} = \PIXCON^\eta\cdot\TR_{\agngvt,\gvt}^{O}$ \hfill 5
		
		
		\item $\TR_{h,\firm} = \lambda^{\TR}_{h,\firm} \cdot\YDF$ \hfill 2		
		
		\item $\TR_{h,\ngo} = \PIXCON^\eta\cdot\TR_{h,\ngo}^{O}$ \hfill 2
		
		\item $\TR_{\ngo,\firm} = \lambda^{\TR}_{\ngo,\firm} \cdot\YDF$  \hfill 1
		
	\end{enumerate}
	
	\redparagraph{Bloc Demande}\hfill \textcolor{Fuchsia}{\textbf{30}}
	\begin{enumerate}[itemsep=2pt,label=\textbf{eq\arabic*}.,resume,ref=\textbf{eq\arabic*}]
		
		\item $\PC_{tr}\cdot\CH_{\tr,h}  = \gamma^{\CH}_{\tr,h} \cdot \CTH_h$   \label{it:CHtrh} \hfill  10
		
		\item $\FBCF = \IT - \sum_{\bns}\PC_{\bns}\cdot\VSTK_{\bns}$ \hfill 1
		
		\item $\PC_{tr}\cdot\INV_{\tr}   = \gamma^{\INV}_{\tr} \cdot \FBCF$   \label{it:INV} \hfill 5 
		
		\item $\DIT_{\tr} = \sum_{j} \DI_{\tr,j}$ \hfill 5
		
		\item $\mu^{D}\cdot Q_{\petr} = \sum_{h} \CH_{\petr,h} + \DIT_{\petr}  + \INV_{\petr}$  \label{it:muDQ}\hfill 1
		
		\item $\MRGCOM_{\bns} = \tmgc_{\bns} \cdot \left(\DD_{\bns} + \IM_{\bns}\right)$ \label{it:MGCOM}  \hfill 4
		
		\item $\MRGTRP_{\bns} = \tmgt_{\bns} \cdot \left( \DD_{\bns} + \IM_{\bns}\right)$  \label{it:MGTRP} \hfill 4
		
	\end{enumerate}
	
	
	\redparagraph{Bloc Commerce international}\hfill \textcolor{Fuchsia}{\textbf{18}}
	\begin{enumerate}[itemsep=2pt,label=\textbf{eq\arabic*}.,resume,  ref=\textbf{eq\arabic*}]
		
		\item $\XS_{\trj} = A_{\trj}^E \cdot \left[\beta_{\trj}^E\cdot\EX_{\trj}^{\rho_{\trj}^E} + \left(1 - \beta_{\trj}^E\right)\cdot \DX_{\trj}^{\rho_{\trj}^E}\right]^{\frac{1}{^{\rho_{\trj}^E}}}$ \hfill 4
		
		\item  $\dfrac{\EX_{\trj}}{\DX_{\trj}} = \left[\left(\dfrac{\PE_{\trj}}{\displaystyle\sum_{\tr}\theta_{\tr,\trj}\cdot\PL_{\trj}}\right)\cdot\left(\dfrac{1-\beta_{\trj}^E}{\beta_{\trj}^E}\right)\right]^{\sigma_{\trj}^E}$, avec $\sigma^E_{\trj} = \dfrac{1}{\rho_{\trj}^E-1}$  \label{it:EX} \hfill 4
		
		\item $\EX_{\petr} = \mu^{E}_{}\cdot Q_{\petr}$ \label{it:EXpetr}\hfill 1
		
		\item $Q_{\trj} = A_{\trj}^M \cdot \bigg[\beta_{\trj}^M\cdot\IM_{\trj}^{-\rho_{\trj}^M} + \left(1 - \beta_{\trj}^M\right)\cdot\DD_{\trj}^{-\rho_{\trj}^M}\bigg]^{-\frac{1}{^{\rho_{\trj}^M}}}$ \label{it:Q} \hfill 4 
		
		
		
		\item  $\dfrac{\IM_{\trj}}{\DD_{\trj}} = \left[\left(\dfrac{\PD_{\trj}}{\PM_{\trj}}\right)\cdot\left(\dfrac{\beta_{\trj}^E}{1-\beta_{\trj}^E}\right)\right]^{\sigma_{\trj}^M}$, avec $\sigma^M_{\trj} = \dfrac{1}{\rho_{\trj}^M+1}$ \label{it:IM}  \hfill 4
		
		\item $\IM_{\petr} = Q_{\petr}$ \hfill 1
		
		
	\end{enumerate}
	
	\redparagraph{Bloc Prix}\hfill \textcolor{Fuchsia}{\textbf{49}} 
	\begin{enumerate}[itemsep=2pt,label=\textbf{eq\arabic*}.,resume,  ref=\textbf{eq\arabic*}]
		
		\item $P_{j}\cdot \XS_{j} = \PVA_{j}\cdot\VA_{j} + \PCI_{j}\cdot\CI_{j}$ \hfill 6
		
		\item $\PCI_j\cdot \CI_j = \displaystyle\sum_{\tr}\PC_{\tr}\cdot\DI_{\tr,j}$ \hfill 6
		
		\item  $\PPD_{\tr} = \PL_{\tr}\cdot\left(1 + \tx_{\tr}\right)$ \label{it:PPD}\hfill 5 
		
		\item $\PD_{\bns} = \PPD_{\bns} + \left(1 + \tx_{\bns}\right)\cdot\left(\tmgc_{\bns} + \tmgt_{\bns}\right)\cdot\PC_{\ser}$ \label{it:PDbns}\hfill 4 
		
		\item $\PD_{\ser} = \PPD_{\ser}$  \label{it:PDser}\hfill 1
		
		\item $\PPM_{\tr} = e\cdot\PWM_{\tr}\cdot\left(1+\tm_{\tr}\right)\cdot\left(1+\tx_{\tr}\right)$\label{it:PPM} \hfill 5  
		
		\item $\PM_{\bns} = \PPM_{\bns} +\left(1 + \tx_{\bns}\right)\cdot\left(\tmgc_{\bns} + \tmgt_{\bns}\right)\cdot\PC_{\ser}$\label{it:PMbns}\hfill 4
		
		\item $\PM_{\ser} = \PPM_{\ser}$ \label{it:PMser}\hfill 1
		
		\item $\PC_{\tr}\cdot Q_{\tr} = \PD_{\tr}\cdot \DD_{\tr} +\PM_{\tr}\cdot\IM_{\tr}$  \label{it:PCbns}  \hfill 5
		
		\item $\PE_{\tr} = e\cdot\PWE_{\tr}$ \hfill 5
		
		\item $P_{\trj}\cdot\XS_{\trj} =  \sum_{\tr}\PL_{\tr}\cdot\DS_{\tr,\trj} + \PE_{\trj}\cdot\EX_{\trj}$ \hfill 4
		
		\item $\PIXPIB = \dfrac{\displaystyle\sum_{j}\PVA_j\cdot\VA_j^O}{\displaystyle\sum_{j}\PVA_j^O\cdot\VA_j^O}$\hfill 1
		
		\item $\PIXCON = \dfrac{\displaystyle\sum_{\tr}\left(\PC_{\tr}\cdot\sum_{h}\CH_{\tr,h}^O\right)}{\displaystyle\sum_{\tr}\left(\PC_{\tr}^O\cdot\sum_{h}\CH_{\tr,h}^O\right)}$ \hfill 1
		
		\item $\PIXINV = \prod_{\tr}\left( \dfrac{\PC_{\tr}}{\PC_{\tr}^O}\right)^{\gamma^{\INV}_{\tr}}$ \hfill 1
		
	\end{enumerate}
	
	\redparagraph{Bloc Équilibre}\hfill \textcolor{Fuchsia}{\textbf{17}} 
	\begin{enumerate}[itemsep=2pt,label=\textbf{eq\arabic*}.,resume,  ref=\textbf{eq\arabic*}]
		
		\item $Q_{\npbns} = \sum_{h} \CH_{\npbns,h} + \DIT_{\npbns} + \INV_{\npbns} +  \VSTK_{\npbns}$  \label{it:Qpetr}\hfill 3
		
		\item $ \PC_{\petr}\cdot Q_{\petr} = \PC_{\petr}\cdot\left(\sum_{h} \CH_{\petr,h} + \DIT_{\petr}  + \INV_{\petr} + \VSTK_{\petr}\right) + \PE_{\petr}\cdot\EX_{\petr}$  \label{it:QDpetr}\hfill 1
		
		
		
		\item $P_{\ntr}\cdot \XS_{\ntr} = \CG_{\ntr}$ \hfill 2
		
		\item $\sum_{\trj}\DS_{\tr,\trj} = \DD_{\tr}$ \hfill 5
		
		\item $\LS = \LD^R+\sum_{j}\LD_j$ \label{it:LS}\hfill 1
		
		\item $\KS_{\trj} = \KD_{\trj}$ \label{it:KS}\hfill 4
		
		\item $\IT  = \sum_{h}\SH_h + \SF + \SG + \SNG + e\cdot\SROW $  \label{it:IT}\hfill 1
		
	\end{enumerate}
	
	\redparagraph{Bloc Produit Intérieur Brut}\hfill \textcolor{Fuchsia}{\textbf{4}} 
	\begin{enumerate}[itemsep=2pt,label=\textbf{eq\arabic*}.,resume,  ref=\textbf{eq\arabic*}]
		
		\item $ \PIB^{\mathit{\mathit{BP}}} = \sum_{j} \PVA_{j}\cdot\VA_{j}$ \hfill 1 
		
		\item $ \PIB^{\mathit{MP}} =  \PIB^{\mathit{BP}} + \sum_{\tr} 	\TIM_{\tr} +\sum_{\tr} \TI_{\tr}$ \hfill 1 
		
		\item $ \PIB^{\mathit{IB}} =  \sum_{j} W\cdot\LD_{j} + \sum_{\trj} R_{\trj}\cdot\KD_{\trj} + \sum_{\tr} \TIM_{\tr} +\sum_{\tr} \TI_{\tr}$ \hfill 1 
		
		\item $\begin{aligned}[t]
			\PIB^{\mathit{FD}} &=  \sum_{\tr}\sum_{h} \PC_{\tr}\cdot\CH_{\tr,h} + \sum_{\ntr} \CG_{\ntr} +  \sum_{\bns}\PC_{\bns}\cdot\left(\INV_{\bns} + \VSTK_{\bns}\right)  \\
			&+ \PC_{\ser}\cdot\INV_{\ser} +   e\cdot\sum_{\tr}^{} \PWE_{\tr}\cdot\EX_{\tr} - e\cdot \sum_{\tr}^{} \PWM_{\tr}\cdot\IM_{\tr}
		\end{aligned}$   \hfill 1
		
	\end{enumerate}
	
	
	\redparagraph{Bloc Variables réels (calculées à partir des indices de prix)}\hfill \textcolor{Fuchsia}{\textbf{7}} 
	\begin{enumerate}[itemsep=2pt,label=\textbf{eq\arabic*}.,resume,  ref=\textbf{eq\arabic*}]
		
		\item $ \CTH_h^{\mathit{reel}} = \dfrac{\CTH_h}{\PIXCON}$ \hfill 2 
		
		\item $ \CG_{\ntr}^{\mathit{reel}} = \dfrac{\CG_{\ntr}}{P_{\ntr}}$ \hfill 2 
		
		\item $ \PIB^{\mathit{BP}_\mathit{reel}} = \dfrac{\PIB^{BP}}{\PIXPIB}$ \hfill 1 
		
		\item $ \PIB^{\mathit{MP}_\mathit{reel}} = \dfrac{\PIB^{MP}}{\PIXCON}$ \hfill 1 
		
		\item $ \FBCF^{\mathit{reel}} = \dfrac{\FBCF}{\PIXINV}$ \hfill 1 		
		
	\end{enumerate}
	
	
	
	
	\redparagraph{Loi de Walras (marché des biens et services)}\hfill \textcolor{Fuchsia}{\textbf{1}} 
	\begin{enumerate}[itemsep=2pt,label=\textbf{eq\arabic*}.,resume,  ref=\textbf{eq\arabic*}]
		
		\item $\LEON  = Q_{\ser}  - \sum_{h} \CH_{\ser,h} - \DIT_{\ser}	 -  \sum_{\bns}\MRGCOM_{\bns} - \sum_{\bns}\MRGTRP_{\bns} - \INV_{\ser}$ \label{it:walras} \hfill 1
		
		\hfill \textcolor{Fuchsia}{$\overline{\textbf{Total (bloc réel): 271}}$}
	\end{enumerate}
	
	
	\redparagraph{Bloc financier}\hfill \textcolor{Fuchsia}{\textbf{26}} 
	\begin{enumerate}[itemsep=2pt,label=\textbf{eqfin\arabic*}.,  ref=\textbf{eqfin\arabic*}]
		
		\item $\TOTG = \CG_{\pub} + \sum_{\agngvt\neq\row} \TR_{\agngvt,\gvt} + e\cdot\TR_{\row,\gvt} + \INVG$\hfill 1
		
		\item $\BRG = \INVG - \SG $\hfill 1
		
		\item $\DEFG = - \BRG $\hfill 1
		
		\item $\INVG = \DINVG + e\cdot\FINVG $\hfill 1
		
		\item $\FINVG = \KTRG + \Delta\NFDG^{\textit{prj}} $\hfill 1
		
		\item $\Delta\NFDG = \Delta\NFDG^{\textit{prg}} + \Delta\NFDG^{\textit{prj}} -\AMORTG$\hfill 1
		
		\item $\BRG = e\cdot\KTRG + e\cdot\Delta\NFDG + \Delta\DCG + \Delta\BG$ \hfill 1
		
		\item $\SP = \SF + \SNG + \sum_{h} \SH_h$ \hfill 1
		
		\item $\Delta\BG = \psi^{\BG}\cdot \SP$ \hfill 1
		\item $\MD = \frac{1}{\vel}\cdot \PIB^{\textit{MP}}$ \hfill 1
		
		\item 	$\MS = \MD$	\hfill 1
		
		\item 	$\Delta\MS  = \ER\cdot \Delta \RES + \Delta\DCG + \Delta\DCP + \RES_0\cdot\Delta\ER$ \hfill 1	
		
		\item 	$\Delta\ER  = \ER -\ER_0$ \hfill 1				
		
		\item 	$\MD  = \MD_0 + \Delta \MD$ \hfill 1				
		
		\item 	$\MS  = \MS_0 + \Delta \MS$ \hfill 1	
		
		\item 	$\DCG  = \DCG_0 + \Delta \DCG$ \hfill 1		
		
		\item 	$\DCP  = \DCP_0 + \Delta \DCP$ \hfill 1												
		
		\item 	$\RES  = \RES_0 + \Delta \RES$ \hfill 1	
		
		\item 	$\Delta \RES = \CAB + \KTRG + \OTHKTR + \FDI^{\textit{NET}} + \Delta \NFDG + \Delta \NFDP $ \hfill 1	
		
		\item 	$\FDI^{\textit{NET}} = \FDI^{\textit{IN}} - \FDI^{\textit{OUT}} $ \hfill 1							
		
		\item 	$\INVP  = \IT - \INVG$ \hfill 1		
		
		\item 	$\Delta\DCP = \psi^{\DCP} \cdot \INVP$ \hfill 1	
		
		\item 	$\NFDG  = \NFDG_0 + \Delta \NFDG$ \hfill 1		
		
		\item 	$\NDDG  = \NDDG_0 + \Delta \DCG+ \Delta \BG$ \hfill 1							
		
		\item 	$\NFDG  = \psi^{\EX} \cdot \sum_{\tr}\PWE_{\tr}\cdot\EX_{\tr}$ \hfill 1		
		
		\item 	$\DEBTG = \ER\cdot\NFDG + \NFDP$ \hfill 1			
	\end{enumerate}
	
	\redparagraph{Loi de Walras (marché financier)}\hfill \textcolor{Fuchsia}{\textbf{1}} 
	\begin{enumerate}[itemsep=2pt,label=\textbf{eqfin\arabic*}.,resume,  ref=\textbf{eqfin\arabic*}]
		\item 	$ \begin{aligned}[t]
			\LEONF & = \SP - \INVP + \Delta\DCP + e \cdot \left(\OTHKTR + \FDI^{\textit{NET}} + \Delta\NFDP\right)\\
			& + \RES_0 \cdot \Delta \ER + (\ER - e) \cdot \Delta\RES - \Delta\MD - \Delta\BG
		\end{aligned}$   \hfill 1
		
		\hfill \textcolor{Fuchsia}{$\overline{\textbf{Total (bloc financier): 27}}$}
		
		\hfill \textcolor{Fuchsia}{$\overline{\textbf{Total équations: 298}}$}
	\end{enumerate}
	
	\subsection{Liste des variables}
	\redparagraph{Variables en volume}\hfill \textcolor{Fuchsia}{\textbf{148}} 
	\begin{description}[leftmargin=!,labelwidth=\widthof{\bfseries AAAAAAAA},parsep=10pt]
		\item[$\CG_{\ntr}^{\mathit{reel}}$] Consommation réelle finale de l'agent $\ntr$ \hfill 2
		\item[$\CH_{\tr,h}$]	Consommation du ménage $h$ en produit $\tr$ 					\hfill 10
		\item[$\CI_j$]      	Consommation intermédiaire totale de la branche $j$ 			\hfill 6
		\item[$\CTH_h^{\mathit{reel}}$]	Consommation réelle finale du ménage  $h$ 				\hfill 2
		\item[$\DD_{\tr}$]		Demande de produit local $\tr$ 									\hfill 5
		\item[$\DI_{\tr,j}$]	Consommation intermédiaire en produit $\tr$ par la branche $j$ 	\hfill 30
		\item[$\DIT_{\tr}$]		Demande intermédiaire totale pour le produit $\tr$ 				\hfill 5
		\item[$\DS_{\tr,\trj}$]	Offre de produit $\tr$ de la branche $\trj$ sur le marché local \hfill 20
		\item[$\DX_{\trj}$]		Production de la branche $\trj$ sur le marché local 			\hfill 4
		\item[$\EX_{\tr}$]		Exportations du produit $\tr$ 		\hfill 5
		\item[$\FBCF^{\mathit{reel}}$]			Formation brute de capital fixe 	\hfill 1
		\item[$\IM_{\tr}$]		Importations du produit $\tr$ 	\hfill 5
		\item[$\INV_{\tr}$]		Demande finale en produit $\tr$ pour fins d'investissement 	\hfill 5
		\item[$\KD_{\trj}$]		Demande de capital de la branche $\trj$ 		\hfill 4
		\item[$\KD^R$]			Demande de capital du Reste du monde 	\hfill 1
		\item[$\KS_{\trj}$]		Offre de capital dans la branche $\trj$ 		\hfill 4
		\item[$\LD_{j}$]		Demande de travail de la branche $j$ 		\hfill 6
		\item[$\LD^R$]			Demande de travail du Reste du monde 	\hfill 1
		\item[$\LS$]			Offre totale de travail  				\hfill 1
		\item[$\MRGCOM_{\bns}$]	Marge commerciale du produit $\bns$ 	\hfill 4
		\item[$\MRGTRP_{\bns}$]	Marge de transport du produit $\bns$ 	\hfill 4
		\item[$\PIB^{\mathit{MP}_\mathit{reel}}$] Produit Intérieur Brut (PIB) aux prix du marche \hfill 1
		\item[$\PIB^{\mathit{BP}_\mathit{reel}}$] PIB aux prix de base 	\hfill 1
		\item[$Q_{\tr}$]		Demande en produit composite $\tr$ 		\hfill 5
		\item[$\VA_j$]			Valeur ajoutée de la branche $j$ 		\hfill 6
		\item[$\VSTK_{\bns}$]	Variation de stocks du produit $\bns$	\hfill 4
		\item[$\XS_j$]			Production de la branche $j$ 			\hfill 6
	\end{description}
	%		
	\redparagraph{Prix}\hfill \textcolor{Fuchsia}{\textbf{74}} 
	\begin{description}[leftmargin=!,labelwidth=\widthof{\bfseries AAAAAAAA},parsep=10pt]
		\item[$e$]			Taux de change nominal 					\hfill 1
		\item[$P_j$]		Prix au producteur de la branche $j$ 	\hfill 6
		\item[$\PC_{\tr}$]	Prix du produit composite $\tr$ 		\hfill 5
		\item[$\PCI_j$]		Indice des prix à la consommation intermédiaire de la branche $j$ 	\hfill 6
		\item[$\PD_{\tr}$]	Prix du produit $tr$ sur le marche local (taxes et marges incluses)	\hfill 5
		\item[$\PE_{\tr}$]	Prix au producteur du produit exporté $\tr$ 	\hfill 5
		\item[$\PIXPIB$]	Déflateur du PIB 								\hfill 1
		\item[$\PIXCON$]    Indice des prix à la consommation 				\hfill 1
		\item[$\PIXINV$]    Indice des prix des investissements				\hfill 1
		\item[$\PL_{\tr}$]	Prix du produit $\tr$ sur le marche local (hors taxes et hors marges)				\hfill 5
		\item[$\PM_{\tr}$]	Prix du produit importe $\tr$ (droits de douane et taxes inclus; marges incluses)	\hfill 5
		\item[$\PPD_{\tr}$]	Prix du produit $\tr$ sur le marche local (taxes incluses; hors marges)				\hfill 5
		\item[$\PPM_{\tr}$]	Prix du produit importe $\tr$ (taxes incluses; hors marges) \hfill 5
		\item[$\PVA_j$]		Prix de la valeur ajoutée de la branche $j$ 				\hfill 6
		\item[$\PWE_{\tr}$]	Prix mondial (en devises) du produit exporté $\tr$ 			\hfill 5
		\item[$\PWM_{\tr}$]	Prix mondial (en devises) du produit importé $\tr$ 			\hfill 5
		\item[$R_{\trj}$]	Taux de rémunération du capital de la branche $\trj$		\hfill 4
		\item[$R^R$]		Taux de rémunération du capital du Reste du monde (en devises) 		\hfill 1
		\item[$W$]			Taux de salaire des travailleurs domestiques				\hfill 1
		\item[$W^R$]		Taux de salaire des travailleurs du Reste du monde (en devises)		\hfill 1
		
	\end{description}
	%		
	\redparagraph{Variables nominales}\hfill \textcolor{Fuchsia}{\textbf{74}} 
	\begin{description}[leftmargin=!,labelwidth=\widthof{\bfseries AAAAAAAA},parsep=10pt]
		
		\item[$\CAB$]	Solde du compte courant (en devises)				\hfill 1
		\item[$\CTH_h$]	Budget de consommation du ménage $h$ 	\hfill 2
		\item[$\FBCF$]	Formation brute de capital fixe 		\hfill 1%4
		\item[$\CG_{\ntr}$]	Budget de consommation finale de l'agent $\ntr$	\hfill 2
		\item[$\IT$]	Investissement total 								\hfill 1
		\item[$\NTXF$]   Recettes non fiscales provenant des Sociétés reçus par l'Administration Centrale \hfill 1
		\item[$\NTXH_h$]  Recettes non fiscales provenant du ménages $h$ \hfill 2
		\item[$\NTXHH_{\agpub}$] Recettes non fiscales provenant des ménages reçus par l'agent public $\agpub$ \hfill 2
		
		\item[$\PIB^{\mathit{BP}}$]	PIB nominal aux prix de base			\hfill 1
		\item[$\PIB^{\mathit{FD}}$]	PIB nominal (approche demande finale)	\hfill 1
		\item[$\PIB^{\mathit{IB}}$]	PIB nominal (approche revenus)			\hfill 1
		\item[$\PIB^{\mathit{MP}}$]	PIB nominal aux prix du marché			\hfill 1%7
		\item[$\SF$]	Épargne des sociétés 					\hfill 1
		\item[$\SG$]	Épargne de l'Administration centrale 	\hfill 1
		\item[$\SH_h$]	Épargne du ménage $h$ 					\hfill 2
		\item[$\SNG$]	Épargne des autres APU et des ISBLSM 	\hfill 1
		\item[$\SROW$]  Épargne du Reste du Monde (en devises) 				\hfill 1%6
		\item[$\TDF$]	Recettes provenant des taxes directes sur les sociétés 				\hfill 1
		\item[$\TDH_h$]	Recettes provenant des taxes directes sur le ménage $h$ 			\hfill 2
		\item[$\TIF$]	Recettes provenant des autres taxes indirectes des sociétés 		\hfill 1
		\item[$\TIMF$]	Recettes provenant des autres tarifs douaniers des sociétés 		\hfill 1
		\item[$\TIM_{\tr}$]	Recettes provenant des tarifs douaniers sur le produit $\tr$ 	\hfill 5
		\item[$\TI_{\tr}$]	Recettes provenant des taxes indirectes sur le produit $\tr$ 	\hfill 5%15
		\item[$\TRA_{\agngvt,\gvt}$]   Transferts de l'Administration centrale vers l'agent non gouvernemental  $\agngvt$			\hfill 5
		\item[$\TRA_{\agd,\row}$]   Transferts du Reste du monde vers l'agent domestique  $\agd$ (en devises) 			\hfill 5
		\item[$\TRA_{h,\firm}$]   Transferts des sociétés vers le ménage    $h$			\hfill 2
		\item[$\TRA_{\ngo,\firm}$]   Transferts des sociétés vers les autres APU et les ISBLSM   			\hfill 1
		\item[$\TRA_{h,\ngo}$]   Transferts des autres APU et des ISBLSM vers le ménage $h$ \hfill 2
		\item[$\YDH_h$]	Revenu disponible du ménage $h$ 	\hfill 2
		\item[$\YDF$]	Revenu disponible des sociétés	 	\hfill 1
		\item[$\YF$]	Revenu total des sociétés 			\hfill 1
		\item[$\YFK$]	Revenu du capital des sociétés 		\hfill 1%5
		\item[$\YFTR$]	Revenu des transferts reçus par les sociétés 	\hfill 1
		\item[$\YG$]	Revenu total de l'Administration centrale 		\hfill 1%2
		\item[$\YGTR$]	Revenu des transferts reçus par l'Administration centrale 	\hfill 1
		\item[$\YGTX$]	Revenu des recettes fiscales de l'Administration centrale 	\hfill 1%2
		\item[$\YGNTX$]	Revenu des recettes non fiscales de l'Administration centrale 	\hfill 1%2
		\item[$\YH_h$]	Revenu total du ménage $h$ 			\hfill 2
		\item[$\YHK_h$]	Revenu du capital du ménage $h$ 	\hfill 2
		\item[$\YHL_h$]	Revenu du travail du ménage $h$ 	\hfill 2%6
		\item[$\YHTR_h$]Revenu des transferts reçus par le ménage $h$ 	\hfill 2
		\item[$\YNG$]	Revenu total des autres APU et des ISBLSM		\hfill 1
		\item[$\YROW$]	Revenu total du Reste du monde	(en devises)				\hfill 1
		\item[$\YROWK$]	Revenu du capital provenant du Reste du monde (en devises)	\hfill 1
		\item[$\YROWL$]	Revenu du travail provenant du Reste du monde (en devises)	\hfill 1%6
		
	\end{description}
	
	\redparagraph{Variables financières}\hfill \textcolor{Fuchsia}{\textbf{34}} 
	\begin{description}[leftmargin=!,labelwidth=\widthof{\bfseries AAAAAAAA},parsep=10pt]
		\item[$\AMORTG$]          Amortissement de la dette publique (en devises)	\hfill 1
		\item[$\BRG$]             Besoin de financement du secteur publique	\hfill 1
		\item[$\DCG$]             Crédits bancaires au secteur publique (Stock)	\hfill 1
		\item[$\DCP$]             Crédits bancaires au secteur privé (Stock)	\hfill 1
		\item[$\DEBTG$]           (Stock de la) dette publique totale	\hfill 1
		\item[$\DEFG$]            Déficit budgétaire	\hfill 1
		\item[$\DINVG$]           Investissement public sur financement intérieur	\hfill 1
		\item[$\Delta\BG$]            Emprunt de du secteur publique auprès du secteur privé	\hfill 1
		\item[$\Delta\DCG$]           Variation des crédits bancaires au secteur publique	\hfill 1
		\item[$\Delta\DCP$]          Variation des crédits bancaires au secteur prive	\hfill 1
		\item[$\Delta\ER$]            Variation du taux de change de fin de période	\hfill 1
		\item[$\Delta\MD$]      Variation de la demande de monnaie	\hfill 1
		\item[$\Delta\MS$]            Variation de l'offre de monnaie	\hfill 1
		\item[$\Delta\NFDG$]          Financement extérieur net total de l'administration centrale (en devises)	\hfill 1
		\item[$\Delta\NFDG^{\mathit{prg}}$]      Financement extérieur net - Prêts programmes (en devises)	\hfill 1
		\item[$\Delta\NFDG^{\mathit{prj}}$]      Financement extérieur net - Prêts projets (en devises)	\hfill 1
		\item[$\Delta\NFDP$]          Financement extérieur net du secteur privé  (en devises)	\hfill 1
		\item[$\Delta\RES$]           Variation des réserves internationales (en devises)	\hfill 1
		\item[$\ER$]              Taux de change de fin de période	\hfill 1
		\item[$\FDI^{\mathit{IN}}$]          Entrées brutes d'IDE (Flux entrants; en devises)	\hfill 1
		\item[$\FDI^{\mathit{NET}}$]         Flux nets d'Investissements Directs Étrangers (IDE; en devises)	\hfill 1
		\item[$\FDI^{\mathit{OUT}}$]         Sorties brutes d'IDE (Flux sortants; en devises)	\hfill 1
		\item[$\FINVG$]           Investissement public sur financement extérieur (en devises)	\hfill 1
		\item[$\INVG$]            Investissement public total (Dépenses en capital)	\hfill 1
		\item[$\INVP$]            Investissement du secteur prive	\hfill 1
		\item[$\KTRG$]            Transferts de capital reçus par le secteur publique (Dons en capital; en devises)	\hfill 1
		\item[$\MD$]              Demande de monnaie (Masse monétaire)	\hfill 1
		\item[$\MD$]              Offre de monnaie	\hfill 1
		\item[$\NDDG$]            (Stock de la) dette publique intérieure	\hfill 1
		\item[$\NFDG$]            (Stock de la) dette publique extérieure (en devises)	\hfill 1
		\item[$\OTHKTR$]          Transferts de capital reçus par le secteur prive (en devises)	\hfill 1
		\item[$\RES$]             Stock de réserves internationales (Avoirs extérieurs nets; en devises)	\hfill 1
		\item[$\SP$]              Epargne du secteur prive	\hfill 1
		\item[$\TOTG$]            Dépenses publiques totales	\hfill 1
	\end{description}
	
	
	
	%		
	\redparagraph{Autres variables}\hfill \textcolor{Fuchsia}{\textbf{2}} 
	\begin{description}[leftmargin=!,labelwidth=\widthof{\bfseries AAAAAAA}]
		\item[$\LEON$]  Offre excédentaire de service sur le marché des biens et services \hfill 1
		
		\item[$\LEONF$] Équilibre sur le marché financier \hfill 1
		
		\hfill \textcolor{Fuchsia}{$\overline{\textbf{Total: 332}}$}
	\end{description}
	%		
	\redparagraph{Variables exogènes (fermeture)}\hfill \textcolor{Fuchsia}{\textbf{34}}\\
	\textcolor{Fuchsia}{\textbf{Bloc réel}}\hfill \textcolor{Fuchsia}{\textbf{26}}
	\begin{description}[leftmargin=!,labelwidth=\widthof{\bfseries AAAAAAAA},parsep=10pt]	
		\item[$\CAB$]			Solde du compte courant  (en devises)			\hfill 1
		\item[$e$]				Taux de change nominal 				\hfill 1
		\item[$\CG_{\pub}$]		Budget de consommation finale de l'Administration centrale	\hfill 1
		\item[$\KD_{\trj}$]		Demande de capital de la branche $\trj$ 		\hfill 4
		\item[$\KD^R$]			Demande de capital du Reste du monde 	\hfill 1
		\item[$\LS$]			Offre totale de travail  	\hfill 1
		\item[$\PL_{\petr}$]	Prix du produit $\petr$ sur le marché local (hors taxes et hors marges)	\hfill 1
		\item[$\PWE_{\tr}$]		Prix mondial (en devises) du produit exporté $\tr$ 		\hfill 5
		\item[$\PWM_{\tr}$]		Prix mondial (en devises) du produit importé $\tr$ 		\hfill 5
		\item[$R^R$]			Taux de rémunération du capital du Reste du monde (en devises) 	\hfill 1
		\item[$\SNG$]			Épargne des autres APU et des ISBLSM 	\hfill 1
		\item[$\VSTK_{\npbns}$]	Variation de stocks du produit $\npbns$	\hfill 3
		\item[$W^R$]			Taux de salaire des travailleurs du Reste du monde (en devises)\hfill 1
	\end{description}
	
	\noindent\textcolor{Fuchsia}{\textbf{Bloc financier}}\hfill \textcolor{Fuchsia}{\textbf{8}}
	\begin{description}[leftmargin=!,labelwidth=\widthof{\bfseries AAAAAAAA},parsep=10pt]	
		\item[$\DINVG$]           Investissement public sur financement intérieur	\hfill 1
		\item[$\KTRG$]            Transferts de capital reçus par le secteur publique (Dons en capital; en devises)	\hfill 1
		\item[$\AMORTG$]          Amortissement de la dette publique extérieure (en devises)	\hfill 1
		\item[$\OTHKTR$]          Transferts de capital reçus par le secteur prive (en devises)	\hfill 1
		\item[$\Delta\NFDG^{\mathit{prg}}$]      Financement extérieur net - Prêts programmes (en devises)	\hfill 1
		\item[$\ER$]              Taux de change de fin de période	\hfill 1
		\item[$\FDI^{\mathit{IN}}$]          Entrées brutes d'IDE (Flux entrants; en devises)	\hfill 1
		\item[$\FDI^{\mathit{OUT}}$]         Sorties brutes d'IDE (Flux sortants; en devises)	\hfill 1
		
		\hfill \textcolor{Fuchsia}{$\overline{\textbf{Total: 34}}$}
	\end{description}
	
	
	\redparagraph{Données d'initialisation}
	\begin{description}[leftmargin=!,labelwidth=\widthof{\bfseries AAAAAAAA},parsep=10pt]    
		\item[$\RES_0$] Stock de réserves internationales (Avoirs extérieurs nets; en devises, fin de l'année précédant la base)
		\item[$\DCG_0$] Crédits bancaires au secteur public (Stock de fin d'année précédente)    
		\item[$\DCP_0$] Crédits bancaires au secteur privé (Stock de fin d'année précédente)    
		\item[$\MD_0$]  Masse monétaire (Stock de fin d'année précédente)
	\end{description}
	
	
	\subsection{Liste des paramètres}
	\begin{description}[leftmargin=!,labelwidth=\widthof{\bfseries AAAAA},parsep=10pt]
		\item[$A_{\tr}$] Paramètre d'échelle (Cobb-Douglas - \textit{}fonction de production)
		
		\item[$A_{\trj}^E$] Paramètre d'échelle (CET -  fonction de production)
		
		\item[$A_{\tr}^M$] Paramètre d'échelle (CES -  produit composite)
		
		\item[$\aij_{\tr,j}$] Coefficient technique (Leontief - consommation intermédiaire)
		
		\item[$\alpha_{\tr}$] Élasticité (Cobb-Douglas - fonction de production)
		
		\item[$\beta_{\tr}^E$] Paramètre de répartition (CET - fonction de production)		
		
		\item[$\beta_{\tr}^M$] Paramètre de répartition (CES - produit composite)
		
		\item[$\eta$]    Élasticité-prix des transferts		
		
		\item[$\gamma^{\CH}_{\tr,h}$] Part du produit $\tr$ dans le budget de consommation du ménage $h$
		
		\item[$\gamma^{\INV}_{\tr}$] Part du produit $\tr$ dans les dépenses totales d'investissement
		
		\item[$\io_{j}$] Coefficient technique (Leontief - consommation intermédiaire totale)
		
		\item[$\lambda_h$] Part du revenu du capital reçue par le ménage $h$
		
		\item[$\lambda^R$] Part du revenu du capital reçue par le Reste du monde	
		
		\item[$\lambda^{\TR}_{\ag,\agj}$] Part du transfert à l'agent $\ag$ dans le revenu de l'agent $\agj$
		
		\item[$\mu^{\CG}$]   Part des autres APU et des ISBLSM dans la dépense de consommation finale publique totale
		
		\item[$\mu^D$] Part de l'offre de produits pétroliers allouée a la demande domestique
		
		\item[$\mu^E$] Part de l'offre de produits pétroliers allouée a la réexportation
		
		\item[$\ntyf$] Propension des Sociétés à verser des recettes non fiscales
		
		\item[$\ntyh_h$] Propension du ménage $h$ à verser des recettes non fiscales
		
		\item[$\lambda^{\TRA}_{\ngo,\firm}$]   Part du revenu disponible des Sociétés qui va aux autres APU et aux ISBLSM
		
		\item[$\lambda^{\TRA}_{h,\firm}$]   Part du revenu disponible des Sociétés qui va au ménage $h$
		
		\item[$\phi_h$] Part du revenu du travail reçue par le ménage $h$
		
		\item[$\psi_h$] Propension moyenne à épargner du ménage $h$
		
		\item[$\psi^{\BG}$] Propension du secteur prive a souscrire aux titres publics
		
		\item[$\psi^{\DCP}$]          Ratio de financement de l'investissement prive par le credit bancaire
		
		\item[$\psi^{\EX}$]             Ratio dette publique extérieure - exportations
		
		\item[$\rho_{\trj}^E$] Paramètre d'élasticité (CET - fonction de production)
		
		\item[$\rho_{\trj}^M$] Paramètre d'élasticité (CES - produit composite)
		
		\item[$\sigma_{\trj}^E$] Élasticité de transformation (CET - fonction de production)
		
		\item[$\sigma_{\trj}^M$] Élasticité de transformation (CES - produit composite)
		
		\item[$\sntx_{\agpub}$] Part des recettes non fiscales versées par les ménages et reçues par l'agent public $\agpub$
		
		\item[$\syfi$] Part des autres impôts indirects sur le revenu des Sociétés 
		
		\item[$\syfm$] Part des autres droits de douanes sur le revenu des Sociétés 
		
		\item[$\theta_{\tr,\trj}$] Part du produit $\tr$ dans la production domestique de la branche $\trj$
		
		\item[$\tm_{\tr}$] Taux de droits de douane sur les importations en produit $\tr$
		
		\item[$\tmgc_{\bns}$] Taux de marge commerciale du produit $\tr$
		
		\item[$\tmgt_{\bns}$] Taux de marge de transport du produit $\tr$
		
		\item[$\tx_{\tr}$] Taux de droits de douane sur les importations en produit $\tr$
		
		\item[$\tyf$] Taux de taxe directe sur le revenu des entreprises 
		
		\item[$\tyh_h$] Taux de taxe directe sur le revenu du ménage $h$
		
		\item[$v_j$] Coefficient (Leontief - valeur ajoutée)
		
		
	\end{description}
	
	
	\subsection{Les éléments de la MCS}
	\begin{description}[leftmargin=!,labelwidth=\widthof{\bfseries OTH}]
		\item[\sffamily F] Facteurs
		\begin{description}[labelindent=1cm,leftmargin=1cm,labelwidth=\widthof{\bfseries OTHIND}]
			\item[\sffamily LD] Travail
			\item[\sffamily KD] Capital
		\end{description}
		\item[\sffamily AG] Agents économiques
		\begin{description}[labelindent=1cm,leftmargin=1cm,labelwidth=\widthof{\bfseries OTHIND}]
			\item[\sffamily HUR] Ménages urbains
			\item[\sffamily HRU] Ménages ruraux
			\item[\sffamily GVT] Administration centrale
			\item[\sffamily FIRM] Sociétés
			\item[\sffamily NGO] Autres APU et ISBLSM
			\item[\sffamily TDM] Taxes directes
			\item[\sffamily TBS] Taxes indirectes et autres taxes sur les biens et services
			\item[\sffamily TCM] Droits de douanes et autres taxes sur le commerce international
			\item[\sffamily NTX] Prélèvements non fiscaux
			\item[\sffamily ROW] Reste du monde
		\end{description}
		\item[\sffamily J] Branches d'activités
		\begin{description}[labelindent=1cm,leftmargin=1cm,labelwidth=\widthof{\bfseries OTHIND}]
			\item[\sffamily AGR] Agriculture
			\item[\sffamily FOOD] Alimentaire, Boisson, Tabac
			\item[\sffamily OTHIND] Autres Industries
			\item[\sffamily SER] Service marchand
			\item[\sffamily ADM] Administration publique
			\item[\sffamily OTHNTR] Autres services publics et autres services non marchands
		\end{description}
		\item[\sffamily I] Produits domestiques
		\begin{description}[labelindent=1cm,leftmargin=1cm,labelwidth=\widthof{\bfseries OTHIND}]
			\item[\sffamily AGR] Agriculture
			\item[\sffamily FOOD] Alimentaire, Boisson, Tabac
			\item[\sffamily OTHIND] Autres Industries
			\item[\sffamily PETR] Produits pétroliers
			\item[\sffamily SER] Service marchand
			\item[\sffamily ADM] Administration publique
			\item[\sffamily OTHNTR] Autres services publics et autres services non marchands
		\end{description}
		\item[\sffamily X] Produits d'exportation
		\begin{description}[labelindent=1cm,leftmargin=1cm,labelwidth=\widthof{\bfseries OTHIND}]
			\item[\sffamily AGR] Agriculture
			\item[\sffamily FOOD] Alimentaire, Boisson, Tabac
			\item[\sffamily OTHIND] Autres Industries
			\item[\sffamily PETR] Produits pétroliers
			\item[\sffamily SER] Service marchand
		\end{description}
		\item[\sffamily M] Marges
		\begin{description}[labelindent=1cm,leftmargin=1cm,labelwidth=\widthof{\bfseries OTHIND}]
			\item[\sffamily MC] Marges de commercialisation
			\item[\sffamily MT] Marges de transport
		\end{description}
		\item[\sffamily OTH] Accumulation
		\begin{description}[labelindent=1cm,leftmargin=1cm,labelwidth=\widthof{\bfseries OTHIND}]
			\item[\sffamily INV] Investissement
			\item[\sffamily VSTK] Variation des stocks
			\item[\sffamily TOT] Total
		\end{description}
	\end{description}
	
	
	\subsection{Calibration du modèle}
	\mbox{ }
	% -------------------------------------------------------------
	% 1. PARAMETRES DE PRODUCTION
	% -------------------------------------------------------------
	\begin{table}[!htbp]
		\centering
		\caption{Paramètres de production}
		\label{tab:param_production}
		\small
		\begin{tabular}{@{}l c D{.}{.}{1.3}D{.}{.}{1.3}D{.}{.}{1.3}D{.}{.}{1.3}D{.}{.}{1.3}D{.}{.}{1.3}D{.}{.}{1.3}@{}}
			\toprule
			Désignation & Notation & \multicolumn{1}{c}{$\agr$} &  \multicolumn{1}{c}{$\food$} &  \multicolumn{1}{c}{$\othind$}  &  \multicolumn{1}{c}{$\petr$} &  \multicolumn{1}{c}{$\ser$} &  \multicolumn{1}{c}{$\pub$} & \multicolumn{1}{c@{}}{$\othntr$} \\
			\midrule
			Paramètre d'échelle & $A_{\trj}$ & 1.917 & 2.000 & 2.000& \multicolumn{1}{c}{$-$}  & 1.147 & \multicolumn{1}{c}{$-$} &  \multicolumn{1}{c}{$-$} \\
			Élasticité de la valeur ajoutée au travail & $\alpha_{\trj}$ & 0.644 & 0.494 & 0.494& \multicolumn{1}{c}{$-$}  & 0.031 & \multicolumn{1}{c}{$-$} & \multicolumn{1}{c}{$-$} \\
			\addlinespace
			\textit{Coefficients techniques}\\
			\addlinespace
			\quad Valeur ajoutée & $v_j$ & 0.696 & 0.281 & 0.485& \multicolumn{1}{c}{$-$}  & 0.600 & 0.856 & 0.986 \\			
			\quad  Coefficients d'input-output & $a_{ij}$ \\
			\qquad $\agr$ & & 0.766 & 0.893 & 0.017& \multicolumn{1}{c}{$-$}  & 0.006 & 0.000 & 0.000 \\
			\qquad  $\food$ & & 0.047 & 0.062 & 0.001& \multicolumn{1}{c}{$-$}  & 0.016 & 0.000 & 0.000 \\
			\qquad  $\othind$ & & 0.145 & 0.028 & 0.599& \multicolumn{1}{c}{$-$}  & 0.286 & 0.186 & 0.186 \\
			\qquad  $\petr$ & & 0.008 & 0.004 & 0.091& \multicolumn{1}{c}{$-$}  & 0.084 & 0.067 & 0.067 \\
			\qquad  $\ser$ & & 0.033 & 0.013 & 0.292& \multicolumn{1}{c}{$-$}  & 0.608 & 0.746 & 0.746 \\
			\quad Consommation intermédiaire totale & $\io_j$ & 0.250 & 0.594 & 0.413& \multicolumn{1}{c}{$-$}  & 0.346 & 0.130 & 0.012 \\
			\addlinespace
			\quad Matrice de production multi-produit & $\theta_{\tr,\trj}$ \\
			\qquad  $\agr$ & & 0.857 & 0.101 & 0.042& 0.000 & 0.000 & \multicolumn{1}{c}{$-$} & \multicolumn{1}{c}{$-$} \\
			\qquad  $\food$ & & 0.000 & 1.000 & 0.000& 0.000 & 0.000 & \multicolumn{1}{c}{$-$} & \multicolumn{1}{c}{$-$} \\
			\qquad  $\othind$ & & 0.000 & 0.000 & 0.960& 0.000 & 0.040 & \multicolumn{1}{c}{$-$} & \multicolumn{1}{c}{$-$} \\
			\qquad  $\ser$ & & 0.000 & 0.001 & 0.003& 0.000 & 0.996 & \multicolumn{1}{c}{$-$} & \multicolumn{1}{c}{$-$} \\
			\bottomrule
		\end{tabular}
		
		
		\smallskip
		\Centering
		\small
		\noindent
		\textbf{Source}: Calculés à partir de la MCS.
		
	\end{table}
	
	\vfill
	
	% -------------------------------------------------------------
	% 2. PARAMETRES DU COMMERCE EXTERIEUR
	% -------------------------------------------------------------
	\begin{table}[!htbp]
		\centering
		\caption{Paramètres du commerce extérieur}
		\label{tab:param_commerce}
		\small
		\begin{tabular}{@{}l c D{.}{.}{1.3}D{.}{.}{1.3}D{.}{.}{1.3}D{.}{.}{1.3}@{}}
			\toprule
			Désignation	& Notation & \multicolumn{1}{c}{$\agr$} & \multicolumn{1}{c}{$\food$} & \multicolumn{1}{c}{$\othind$} & \multicolumn{1}{c}{$\ser$} \\
			\midrule
			Paramètre d'échelle Armington (Importations) & $A^M_{\trj}$ & 1.056 & 1.597 & 1.844 & 1.328 \\[1.5ex]
			Paramètre d'échelle CET (Exportations) & $A^E_{\trj}$ & 6.914 & 2.602 & 2.049 & 3.729 \\[1.5ex]
			Paramètre d'élasticité CES & $\rho^M_{\trj}$ & 0.267 & 0.267 & 0.267 & 0.267 \\[1.5ex]
			Paramètre d'élasticité CET & $\rho^E_{\trj}$ & 3.062 & 3.062 & 3.062 & 3.062 \\[1.5ex]
			Élasticité de substitution [*]& $\sigma^M_{\trj}$ & 0.789 & 0.789 & 0.789 & 0.789 \\[1.5ex]
			Élasticité de transformation [*] & $\sigma^E_{\trj}$ & 0.485 & 0.485 & 0.485 & 0.485 \\[1.5ex]
			Part de répartition des importations & $\beta^M_{\trj}$ & 0.005 & 0.150 & 0.722 & 0.061 \\[1.5ex]
			Part de répartition des exportations & $\beta^E_{\trj}$ & 0.997 & 0.902 & 0.346 & 0.976 \\
			\bottomrule
		\end{tabular}
		
		\smallskip
		\Centering
		\small
		\noindent
		\textbf{Source}: Calculés à partir de la MCS; [*] \cite{Devarajanetal2023}.
		
	\end{table}
	
	
	
	% -------------------------------------------------------------
	% 3. FISCALITE. MARGES ET DEMANDE FINALE
	% -------------------------------------------------------------
	\begin{table}[!htbp]
		\centering
		\caption{Fiscalité. marges et structure de la demande finale}
		\label{tab:param_fiscalite_demande}
		\small
		\begin{tabular}{@{}l c D{.}{.}{1.3}D{.}{.}{1.3}D{.}{.}{1.3}D{.}{.}{1.3}D{.}{.}{1.3}@{}}
			\toprule
			Désignation & Notation & \multicolumn{1}{c}{$\agr$} & \multicolumn{1}{c}{$\food$} & \multicolumn{1}{c}{$\othind$} & \multicolumn{1}{c}{$\petr$} & \multicolumn{1}{c}{$\ser$} \\
			\midrule
			\textit{A. Fiscalité et marges}\\
			\quad Taux de droits de douane & $\tm_{\tr}$ & 0.133 & 0.135 & 0.132 & 0.268 & 0.000 \\
			\quad Taux de taxation indirecte & $\tx_{\tr}$ & 0.005 & 0.037 & 0.015 & 0.000 & 0.011 \\
			\quad Taux de marge commerciale & $\tmgc_{\bns}$ & 0.156 & 0.151 & 0.152 & 0.172 & \multicolumn{1}{c}{$-$} \\
			\quad Taux de marge de transport & $\tmgt_{\bns}$ & 0.038 & 0.035 & 0.032 & 0.269 & \multicolumn{1}{c}{$-$} \\
			\addlinespace
			\textit{B. Coefficients budgétaires}\\
			\quad Propension marginale à consommer (Ménages urbains) & $\gamma_{\hur}^{\CH}$ & 0.253 & 0.388 & 0.115 & 0.001 & 0.243 \\[1.5ex]
			\quad Propension marginale à consommer (Ménages ruraux) & $\gamma_{\hru}^{\CH}$ & 0.239 & 0.424 & 0.108 & 0.001 & 0.229 \\[1.5ex]
			\quad Structure par produit de l'investissement & $\gamma_{\tr}^{\INV}$ & 0.000 & 0.000 & 0.256 & 0.000 & 0.744 \\
			\bottomrule
		\end{tabular}
		
		\smallskip
		\Centering
		\small
		\noindent
		\textbf{Source}: Calculés à partir de la MCS.
		
	\end{table}
	
	
	
	\begin{table}[!htbp]
		\centering
		\small
		\caption{Paramètres socio-économiques}
		\label{tab:param_socio_transferts}
		\begin{tabular}{@{} l c D{.}{.}{1.3} D{.}{.}{1.3} D{.}{.}{1.3} D{.}{.}{1.3} @{}}
			\toprule
			Désignation & Notation & \multicolumn{1}{c}{$\hur$} & \multicolumn{1}{c}{$\hru$} & \multicolumn{1}{c}{$\firm$} & \multicolumn{1}{c}{$\row$} \\
			\midrule
			\multicolumn{6}{@{}l}{\textit{A. Répartition des revenus des facteurs}} \\[0.5ex]
			\quad Part du revenu du travail & $\phi_h$ & 0.541 & 0.454 & \multicolumn{1}{c}{---} & \multicolumn{1}{c}{---} \\
			\quad Part du revenu du capital & $\lambda_h / \lambda^R$ & 0.038 & 0.251 & \multicolumn{1}{c}{---} & 0.057 \\[1.5ex]
			
			\multicolumn{6}{@{}l}{\textit{B. Fiscalité directe et prélèvements}} \\[0.5ex]
			\quad Taux de taxe directe & $\tyh_h / \tyf$ & 0.023 & 0.008 & 0.035 & \multicolumn{1}{c}{---} \\
			\quad Propension à verser des recettes non fiscales & $\ntyh_h / \ntyf$ & 0.011 & 0.011 & 0.005 & \multicolumn{1}{c}{---} \\
			\quad Taux des autres taxes sur les biens et services & $\syfi$ & \multicolumn{1}{c}{---} & \multicolumn{1}{c}{---} & 0.008 & \multicolumn{1}{c}{---} \\
			\quad Taux des autres droits de douane & $\syfm$ & \multicolumn{1}{c}{---} & \multicolumn{1}{c}{---} & 0.027 & \multicolumn{1}{c}{---} \\[1.5ex]
			
			\multicolumn{6}{@{}l}{\textit{D. Comportement d'épargne}} \\[0.5ex]
			\quad Propension à épargner & $\psi_h$ & 0.187 & 0.176 & \multicolumn{1}{c}{---} & \multicolumn{1}{c}{---} \\
			\bottomrule
		\end{tabular}
		
		
		\smallskip
		\Centering
		\small
		\noindent
		\textbf{Source}: Calculés à partir de la MCS.
	\end{table}
	
	
	\begin{table}[!htbp]
		\centering
		\small
		\caption{Matrice des paramètres et valeurs initiales des transferts}
		\label{tab:matrice_transferts}
		\begin{tabular}{@{} l c D{.}{.}{4.3} D{.}{.}{4.3} D{.}{.}{4.3} D{.}{.}{4.3} @{}}
			\toprule
			& & \multicolumn{4}{c}{Émetteurs des transferts ($\ag$)} \\
			\cmidrule(lr){3-6}
			Bénéficiaires & Notation & \multicolumn{1}{c}{$\firm$} & \multicolumn{1}{c}{$\gvt$} & \multicolumn{1}{c}{$\ngo$} & \multicolumn{1}{c}{$\row$} \\
			\midrule
			\multicolumn{6}{@{}l}{\textit{A. Propensions à partir du revenu disponible}} \\[1.5ex]
			\quad Ménages urbains ($\hur$)       & $\lambda^{\TRA}_{\hur, \firm}$  & 0.039 & \multicolumn{1}{c}{$-$} & \multicolumn{1}{c}{$-$} & \multicolumn{1}{c}{$-$} \\[1.5ex]
			\quad Ménages ruraux ($\hru$)        & $\lambda^{\TRA}_{\hru, \firm}$  & 0.663 & \multicolumn{1}{c}{$-$} & \multicolumn{1}{c}{$-$} & \multicolumn{1}{c}{$-$} \\[1.5ex]
			\quad Autres APU et ISBLSM ($\ngo$)  & $\lambda^{\TRA}_{\ngo, \firm}$  & 0.221 & \multicolumn{1}{c}{$-$} & \multicolumn{1}{c}{$-$} & \multicolumn{1}{c}{$-$} \\[1.5ex]
			
			\multicolumn{6}{@{}l}{\textit{B. Transferts forfaitaires initiaux (en milliards d'Ariary, année de base)}} \\[1.5ex]
			\quad Ménages urbains ($\hur$)       & $\TRA_{\hur, \ag}^O$ & \multicolumn{1}{c}{$-$} &  257.537 & 22.968 &  524.466 \\[1.5ex]
			\quad Ménages ruraux ($\hru$)        & $\TRA_{\hru, \ag}^O$ & \multicolumn{1}{c}{$-$} &  354.500 & 32.507 &  742.293 \\[1.5ex]
			\quad Sociétés ($\firm$)             & $\TRA_{\firm, \ag}^O$ & \multicolumn{1}{c}{$-$} & 1292.532 & \multicolumn{1}{c}{$-$} & 1189.735 \\[1.5ex]
			\quad Administration centrale ($\gvt$)           & $\TRA_{\gvt, \ag}^O$ & \multicolumn{1}{c}{$-$} & \multicolumn{1}{c}{$-$} & \multicolumn{1}{c}{$-$} &  365.104 \\[1.5ex]
			\quad Autres APU et ISBLSM ($\ngo$)  & $\TRA_{\ngo, \ag}^O$ & \multicolumn{1}{c}{$-$} &  754.849 & \multicolumn{1}{c}{$-$} &   66.891 \\[1.5ex]
			\quad Reste du monde ($\row$)        & $\TRA_{\row, \ag}^O$ & \multicolumn{1}{c}{$-$} &  106.819 & \multicolumn{1}{c}{$-$} & \multicolumn{1}{c}{$-$} \\
			\bottomrule
		\end{tabular}
		
		\smallskip
		\Centering
		\small
		\noindent
		\textbf{Source}: Calculés à partir de la MCS.
		
	\end{table}
	
	\begin{table}[!htbp]
		\centering
		\small
		\caption{Autres paramètres}
		\label{tab:parametres_simplifie}
		\begin{tabular}{@{} l c D{.}{.}{3.3} @{}}
			\toprule
			Désignation & Notation & \multicolumn{1}{c}{Valeur} \\
			\midrule
			Élasticité d'indexation des transferts[*]  & $\eta$ & 1.000\\[0.5ex]
			Part des autres APU et ISBLSM dans la dép.~de cons.~finale publique & $m_G$ & 0.717 \\[1.5ex]
			
			\multicolumn{3}{@{}l}{\textit{Paramètres du secteur des produits pétroliers :}} \\[0.5ex]
			\quad Part de l'offre allouée à la demande domestique & $\mu^D$ & 0.639 \\
			\quad Part de l'offre allouée à la réexportation & $\mu^E$ & 0.119 \\[1.5ex]
			
			\multicolumn{3}{@{}l}{\textit{Clé de répartition des ressources non fiscales des ménages :}} \\[0.5ex]
			\quad Part captée par l'Administration centrale & $\sntx_{\gvt}$ & 0.054 \\
			\quad Part captée par les autres APU et les ISBLSM & $\sntx_{\ngo}$ & 0.946 \\
			
			\bottomrule
		\end{tabular}
		
		
		\smallskip
		\Centering
		\small
		\noindent
		\textbf{Source}: Calculés à partir de la MCS. [*] \cite{Decaluweetal2013}
		
	\end{table}
	%	\section{MCS}
	%	
	%
	%	
	
	%	
	%	\pagestyle{plain} % Réactive les numéros de page standard en bas au centre
	%	\sffamily
	%	%	\includegraphics{Exter_MadaJ_cropped}
	%	%\incgraph[a3paper, landscape,	overlay page number at bottom][width=0.2\paperheight]{MIDISAM}
	
	\newpage
	\begin{thebibliography}{99}	
		\setlength{\itemsep}{0.25cm}
		\bibitem[Decaluwé et al., 2001]{Decaluweetal2001}
		Bernard Decaluwé, André Martens, et Luc Savard (2001), \textit{La Politique Économique du Développement et les Modèles d'Équilibre Général Calculable : Une Introduction à l'Application de l'Analyse Mésoéconomique aux Pays en Développement}, Les Presses de l'Université de Montréal.
		
		
		\bibitem[Decaluwé et al., 2013]{Decaluweetal2013}
		Bernard Decaluwé, André Lemelin, Véronique Robichaud, and Hélène Maisonnave (2013), \href{https://www.pep-net.org/research-resources/cge-models}{\textit{{PEP-1-1}: The PEP Standard Single-country, Static CGE Model}}, Partnership for Economic Policy (PEP).
		
		\bibitem[Devarajan et al., 2023]{Devarajanetal2023}
		Shantayanan Devarajan, Delfin S.~Go, and Sherman Robinson (2023), ``\href{http://hdl.handle.net/10986/39981}{Trade elasticities in aggregate models: Estimates for 191 countries}'', \textit{Policy Research Working Papers}, \textbf{10490}, World Bank. 
		
		
		\bibitem[Dissou, 1998]{Dissou1998}
		Yazid Dissou (1998), \textit{Modèle d'Équilibre Général de Madagascar (MADCGE): Documentation et Référence}, Institut National de la Statistique (INSTAT), Cornell Food and Nutrition Policy Program.
		
		\bibitem[Dissou, 2001]{Dissou2001}
		Yazid Dissou (2001), \textit{Atelier de Formation sur le Modèle d'Équilibre géneral de Madagascar 1999}, Institut National de la Statistique (INSTAT), Cornell Food and Nutrition Policy Program.
		
		\bibitem[Dissou et al., 1998]{Dissouetal1998}
		Yazid Dissou, Steven Haggblade, Hery Andriamasy, Francis Andrianarison, Michel Botomazava, Charline Raholimalala, Dodo Ralantoarilolona, Christian Rasolomanana, Rajaonson Ndriamahavory, Joé Rabeantoandro, Jean Razafindravonona, Harivelo Rajemison, Josée Randriamamonjy, Jean-Gabriel Randrianarison, Andriamampianina Rakotomalala, Eric Rakotomalala, Tiaray Razafimanantena, Jean Ratsimanandresy, Jean Pierre Randrianasolo et Adeline Razaiarinivo (1998), ``\href{https://www.documentation.ird.fr/hor/fdi:010019178}{Politique fiscale à Madagascar : options et impacts}'', in Mireille Razafindrakoto et François Roubaud (eds.), \textit{Perspectives Macroéconomiques et Politiques Publiques : La Question Fiscale, Économie de Madagascar}, no.~\textbf{3}, pp.~105--132. 
		
		
		\bibitem[Dorosh et al., 2003]{Doroshetal2003}
		Paul Dorosh, Steven Haggblade, Christen Lungren, Tiaray Razafimanantena, and Zaza Randriamiarana (2003), ``\href{https://www.cfnpp.cornell.edu/images/wp144en.pdf}{Economic Motors for Poverty Reduction in Madagascar}'', \textit{CFNPP Working Paper}, \textbf{144}, Cornell Food And Nutrition Policy Program.
		
		\bibitem[Ramaharo, 2022]{Ramaharo2022} Franck Ramaharo (2022), ``\href{https://mpra.ub.uni-muenchen.de/id/eprint/115226}{The 2019 Merged model for Madagascar}'', MPRA Paper \textbf{115226}, University Library of Munich, Germany. 
	\end{thebibliography}	
	
	
	%		\clearpage
	%		\KOMAoptions{paper=A2, paper=portrait, DIV=last} % Switch to A3 Landscape
	%		\newgeometry{
		%				%		top=2cm,       % Marge haute
		%				%		bottom=4cm,    % Marge basse
		%				%		left=2.5cm,      % Marge gauche
		%				%		right=3cm      % Marge droite
		%				scale=.825
		%			}
	%		
	%		\input{MCS-2019}
	%		
	%		\clearpage
	%		\KOMAoptions{paper=A4, paper=portrait, DIV=last} % Switch back to A4
	%		\restoregeometry % Restore previous margins
\end{document}
